\documentclass[12pt,a4paper]{article}

% ── Packages ──────────────────────────────────────────────────────────
\usepackage[utf8]{inputenc}
\usepackage[T1]{fontenc}
\usepackage{lmodern}
\usepackage{amsmath,amssymb,amsthm,mathtools}
\usepackage{bm}
\usepackage{booktabs}
\usepackage{natbib}
\usepackage{geometry}
\usepackage{hyperref}
\usepackage{cleveref}
\usepackage{algorithm}
\usepackage{algpseudocode}
\usepackage{graphicx}
\usepackage{enumitem}
\usepackage{xcolor}
\usepackage{setspace}

\geometry{margin=1in}
\onehalfspacing

% ── Theorem environments ──────────────────────────────────────────────
\newtheorem{theorem}{Theorem}[section]
\newtheorem{proposition}[theorem]{Proposition}
\newtheorem{lemma}[theorem]{Lemma}
\newtheorem{corollary}[theorem]{Corollary}
\newtheorem{definition}[theorem]{Definition}
\newtheorem{remark}[theorem]{Remark}
\newtheorem{assumption}{Assumption}

% ── Notation shortcuts ────────────────────────────────────────────────
\newcommand{\R}{\mathbb{R}}
\newcommand{\E}{\mathbb{E}}
\newcommand{\Var}{\mathrm{Var}}
\newcommand{\Cov}{\mathrm{Cov}}
\newcommand{\tr}{\mathrm{tr}}
\newcommand{\diag}{\mathrm{diag}}
\newcommand{\vect}{\mathrm{vec}}
\newcommand{\kron}{\otimes}
\newcommand{\xf}{{x^{f}}}
\newcommand{\xs}{{x^{s}}}
\newcommand{\gx}{g_x}
\newcommand{\hx}{h_x}
\newcommand{\Gxx}{G_{xx}}
\newcommand{\Hxx}{H_{xx}}
\newcommand{\gss}{{g_{\sigma\sigma}}}
\newcommand{\hsig}{{h_{\sigma\sigma}}}
\newcommand{\cO}{\mathcal{O}}
\DeclareMathOperator{\specrad}{sp}
\DeclareMathOperator{\chol}{chol}

% ══════════════════════════════════════════════════════════════════════
\title{\Large\textbf{A Pruned Square-Root Unscented Kalman Filter for\\
Bayesian Estimation of Second-Order DSGE Models}}

\author{%
[Authors]\thanks{%
We thank [acknowledgments]. All errors are our own.
Correspondence: [email].}
}
\date{\today}

\begin{document}
\maketitle

\begin{abstract}
We develop a novel filtering methodology---the Pruned Square-Root
Unscented Kalman Filter (Pruned SR-UKF)---for likelihood evaluation
in Dynamic Stochastic General Equilibrium (DSGE) models solved with
second-order perturbation methods. The standard approach augments the
state vector to include both first-order states, their Kronecker
products, and second-order corrections, yielding an
$\cO(n_x + n_x^2 + n_x)$-dimensional linear state space amenable to
the Kalman filter. This augmentation is computationally infeasible
for models with more than approximately 15 state variables, as the
per-step cost scales as $\cO((n_x + n_x^2)^3)$. Our method instead
applies the unscented Kalman filter to the $n_x$-dimensional
first-order state alone, propagates the second-order correction
deterministically from the filtered posterior mean, and handles the
quadratic observation nonlinearity through sigma-point quadrature.
We prove that (i)~the sigma-point approximation captures the
conditional mean of the quadratic observation equation exactly and
its covariance to second order; (ii)~the approximation error in the
second-order state correction is $\cO(\sigma^2)$, where $\sigma$
denotes the shock standard deviation; and (iii)~the per-step cost
is $\cO(n_x^3)$, yielding a speedup of $\cO(n_x^3)$ over the
augmented filter. We embed the filter in a Hamiltonian Monte Carlo
(HMC) posterior sampler with implicit adjoint gradients through the
perturbation solver and validate the methodology on the
\citet{schmitt2003closing} small open economy model, demonstrating
that the Pruned SR-UKF produces posterior estimates statistically
indistinguishable from the augmented Kalman filter at a fraction of
the computational cost.

\medskip
\noindent\textbf{Keywords:}
DSGE models, second-order perturbation, pruning, unscented Kalman filter,
square-root filter, Bayesian estimation, Hamiltonian Monte Carlo.

\medskip
\noindent\textbf{JEL Classification:} C11, C32, C63, E32.
\end{abstract}

\newpage
\tableofcontents
\newpage

% ══════════════════════════════════════════════════════════════════════
%  SECTION 1 — INTRODUCTION
% ══════════════════════════════════════════════════════════════════════
\section{Introduction}\label{sec:intro}

Second-order perturbation methods have become indispensable in
macroeconomics. By retaining quadratic terms in the Taylor expansion
of the model's policy functions around the deterministic steady state,
they capture phenomena that first-order (linear) approximations miss
entirely: risk premia, precautionary savings, welfare costs of
business cycles, and asymmetric impulse responses
\citep{schmitt2004solving, kim2008calculating}. Yet Bayesian estimation
of DSGE models solved to second order remains computationally
demanding, and the difficulty grows rapidly with model size.

The central bottleneck is likelihood evaluation. The pruned
state-space system of \citet{andreasen2018pruned} decomposes the
approximate solution into a first-order stochastic component~$\xf_t$
and a second-order deterministic correction~$\xs_t$, then stacks
both together with $\vect(\xf_t {\xf_t}^\top)$ to form an augmented
state of dimension $n_x + n_x^2 + n_x$. This augmented system is
conditionally linear-Gaussian, so the standard Kalman filter applies.
While analytically elegant, the per-step cost of the Kalman filter
scales as $\cO((2n_x + n_x^2)^3)$, which for a model with $n_x = 30$
state variables implies an augmented dimension of~960 and roughly
$8.8 \times 10^8$ floating-point operations per time step---repeated
for every likelihood evaluation inside a Markov Chain Monte Carlo
(MCMC) sampler. This computational barrier has confined second-order
Bayesian estimation to small-scale models in practice.

\paragraph{Contribution.}
We propose the \emph{Pruned Square-Root Unscented Kalman Filter}
(Pruned SR-UKF), a filtering methodology that operates on the
$n_x$-dimensional first-order state~$\xf_t$ alone, while correctly
incorporating second-order effects through two mechanisms:
\begin{enumerate}[nosep]
\item The \emph{observation equation}
$y_t = \gx(\xf_t + \xs_t) + \tfrac{1}{2}\Gxx(\xf_t \kron \xf_t) + \tfrac{1}{2}\gss + \varepsilon_t$
is quadratic in~$\xf_t$. We handle this nonlinearity with the
unscented transform (UT), which propagates $2n_x + 1$ sigma points
through the observation function and recovers the conditional mean
\emph{exactly} and the conditional covariance to second order---because
the observation is a polynomial of degree two in~$\xf_t$.
\item The \emph{second-order correction} $\xs_t$ is propagated
deterministically from the filtered posterior mean of~$\xf_t$. We
characterize the resulting approximation error and show it is
$\cO(\sigma^2)$, where $\sigma$ is the shock standard deviation.
\end{enumerate}
The per-step cost of our filter is $\cO(n_x^3)$---a speedup of
$\cO(n_x^3)$ relative to the augmented Kalman filter. For the
$n_x = 7$ \citet{schmitt2003closing} model this represents a factor
of approximately~343; for a medium-scale model with $n_x = 30$ the
speedup exceeds $10^6$.

We embed the Pruned SR-UKF in a complete Bayesian estimation pipeline:
maximum a~posteriori (MAP) initialization via L-BFGS, Hessian-based
mass-matrix preconditioning with eigenvalue capping, and Hamiltonian
Monte Carlo (HMC) sampling with implicit adjoint gradients through the
perturbation solver \citep{childers2024differentiable}. The entire
pipeline is implemented in TensorFlow with automatic differentiation
through the filter recursion, enabling gradient-based posterior
exploration without finite-difference approximations.

\paragraph{Related literature.}
Our work bridges four strands of the literature: second-order
perturbation methods
\citep{schmitt2004solving, kim2008calculating, andreasen2018pruned},
nonlinear filtering for DSGE models
\citep{fernandez2007estimating, kollmann2014tractable, ivashchenko2014dsge, andreasen2012non},
square-root sigma-point filtering
\citep{julier1997new, wan2000unscented, van2001square},
and gradient-based Bayesian DSGE estimation
\citep{childers2024differentiable}. We defer a detailed review to
\Cref{sec:literature}.

\paragraph{Outline.}
\Cref{sec:literature} reviews the relevant literature.
\Cref{sec:model} presents the model framework, including the
second-order perturbation solution and the pruned state-space system.
\Cref{sec:method} develops the Pruned SR-UKF algorithm in detail.
\Cref{sec:theory} states and proves the theoretical properties.
\Cref{sec:implementation} discusses implementation and numerical
considerations.
\Cref{sec:challenges} addresses potential objections with
mathematical rebuttals.
\Cref{sec:application} applies the methodology to the
\citet{schmitt2003closing} small open economy model.
\Cref{sec:montecarlo} presents Monte Carlo evidence.
\Cref{sec:conclusion} concludes.


% ══════════════════════════════════════════════════════════════════════
%  SECTION 2 — LITERATURE REVIEW
% ══════════════════════════════════════════════════════════════════════
\section{Literature Review}\label{sec:literature}

\subsection{Second-Order Perturbation Methods and Pruning}

The perturbation approach to solving DSGE models expands the policy
functions in a Taylor series around the deterministic steady state.
\citet{schmitt2004solving} derive a general second-order perturbation
method by reducing the quadratic terms of the equilibrium conditions
to a generalized Sylvester equation. Their method applies to any
rational-expectations model written in the canonical form
$\E_t[H(y_{t+1}, y_t, x_{t+1}, x_t)] = 0$.

A well-known difficulty with second-order (and higher) approximations
is that direct simulation of the full nonlinear law of motion can
produce explosive sample paths, even when the underlying linear
solution is stable. \citet{kim2008calculating} address this by
introducing a ``pruning'' scheme that decomposes the state into a
first-order stochastic component and a second-order correction, and
crucially uses only the first-order component in the quadratic terms.
This prevents the explosive feedback loop while preserving second-order
accuracy.

\citet{andreasen2018pruned} provide a rigorous analysis of the pruned
state-space system for approximations up to third order. They establish
stationarity conditions, derive closed-form expressions for
unconditional first and second moments, and show how the pruned system
can be cast as a linear state-space model suitable for the Kalman
filter---at the cost of augmenting the state vector to include the
Kronecker products of the first-order states.

\subsection{Nonlinear Filtering for DSGE Estimation}

\citet{fernandez2007estimating} pioneer likelihood-based Bayesian
estimation of nonlinear DSGE models using the particle filter
(sequential Monte Carlo). The particle filter provides asymptotically
exact likelihood estimates but requires $\cO(N_p)$ model evaluations
per time step, where $N_p$ is the number of particles---typically
$10^3$ to $10^5$ for acceptable variance.

\citet{kollmann2014tractable} develops a deterministic alternative
by applying the standard Kalman filter to the pruned augmented state
space of \citet{kim2008calculating}. This avoids stochastic simulation
entirely and produces competitive accuracy relative to particle
filters, but inherits the $\cO((n_x + n_x^2)^3)$ cost that limits
scalability.

\citet{ivashchenko2014dsge} compares several nonlinear Kalman filters
for second-order DSGE estimation, including the Unscented Kalman
Filter (UKF), the Central Difference Kalman Filter (CDKF), and the
Quadratic Kalman Filter (QKF). He finds the QKF superior in
parameter-estimation accuracy but does not employ pruning, instead
applying the filters to the full nonlinear dynamics.

\citet{andreasen2012non} applies the CDKF to second- and third-order
DSGE solutions, demonstrating good performance with Gaussian,
Laplace-distributed, and stochastic volatility shocks, again without
the pruning decomposition.

\citet{monfort2015quadratic} develop the Quadratic Kalman Filter for
general state-space models with quadratic measurement and transition
equations, achieving up to 70\% RMSE improvements over the UKF and
extended Kalman filter. However, the QKF requires the state-space to
be written in a specific quadratic form that may not map directly
to the pruned decomposition.

\subsubsection{Why Ivashchenko (2014) and Andreasen (2012) Could Avoid
Pruning---and Why We Cannot}\label{sec:why-pruning}

A natural question arises: if \citet{ivashchenko2014dsge} successfully
applied the UKF, and \citet{andreasen2012non} the CDKF, to
second-order DSGE solutions \emph{without} pruning, why does our
method require it? The answer involves a fundamental distinction
in what these approaches ask the filter to do, and the consequences
for stability, theoretical guarantees, and scalability to full
Bayesian estimation.

\paragraph{The unpruned approach: sigma points through the full
nonlinear dynamics.}
Both Ivashchenko and Andreasen treat the second-order DSGE solution
as a \emph{generic nonlinear state-space model}:
\begin{align}
x_{t+1} &= \hx\,x_t
  + \tfrac{1}{2}\Hxx\,(x_t \kron x_t)
  + \tfrac{1}{2}\hsig + \eta\,\varepsilon_{t+1},
  \label{eq:unpruned-transition} \\
y_t &= \gx\,x_t
  + \tfrac{1}{2}\Gxx\,(x_t \kron x_t)
  + \tfrac{1}{2}\gss + \varepsilon_t^y.
  \label{eq:unpruned-obs}
\end{align}
The UKF (or CDKF) is then applied as a black-box nonlinear filter:
sigma points are propagated through \emph{both} the transition
equation~\eqref{eq:unpruned-transition} and the observation
equation~\eqref{eq:unpruned-obs}. The filtered state is the full
$x_t$---there is no decomposition into first- and second-order
components.

This approach works in practice for several reasons:
\begin{enumerate}[nosep]
\item \textbf{Localized sigma points.}
The UKF/CDKF sigma points are constructed around the current
filtered mean $\hat{x}_{t|t}$ with spread determined by the
filtered covariance $P_{t|t}$. For typical DSGE calibrations with
small shocks ($\sigma_\varepsilon \sim 0.01$--$0.1$), the sigma
points remain close to the steady state, where the quadratic
terms $\tfrac{1}{2}\Hxx(x_t \kron x_t)$ are small relative to the
linear term $\hx\,x_t$. The quadratic feedback that causes
explosions in \emph{open-loop simulation} is suppressed by the
\emph{closed-loop} nature of filtering: each observation corrects
the state estimate back toward the true state, preventing the
accumulation of quadratic drift.

\item \textbf{Point estimation context.}
Ivashchenko uses quasi-maximum likelihood (QML) at fixed parameter
values; Andreasen likewise uses QML. In both cases, the filter
runs at a single parameter configuration (or along a fixed
optimization path). If the filter diverges for a particular
parameter draw, the likelihood is simply $-\infty$ and that
parameter is rejected. Neither author embeds the filter in an
MCMC sampler that must evaluate the likelihood at thousands of
arbitrary parameter configurations.

\item \textbf{Small models.}
Both papers use small DSGE models ($n_x \leq 10$), where the
quadratic terms are moderate and the probability of filter
divergence at reasonable parameter values is low.
\end{enumerate}

\paragraph{Why these reasons are insufficient for Bayesian MCMC
estimation.}
In an MCMC setting---especially with gradient-based samplers like
HMC---the filter must evaluate the likelihood at a wide range of
parameter configurations, including regions where:
\begin{enumerate}[nosep]
\item $\specrad(\hx)$ is close to~1, amplifying the quadratic feedback;
\item $\|\Hxx\|$ is large relative to the shock variance,
making the quadratic terms dominant;
\item the model is near the boundary of Blanchard--Kahn determinacy,
where the dynamics are ill-conditioned.
\end{enumerate}
In these regions, the unpruned sigma-point propagation through
\eqref{eq:unpruned-transition} can produce sigma points that lie far
from the steady state. These outlier sigma points generate large
$x \kron x$ terms, which feed back into the next predict step,
creating the explosive positive-feedback loop that
\citet{kim2008calculating} identified. Even if such explosions are
rare, a single divergent likelihood evaluation during leapfrog
integration can derail an entire HMC trajectory, destroying the
sampler's efficiency.

Moreover, the unpruned system \eqref{eq:unpruned-transition} is
\emph{not stationary}: \citet{andreasen2018pruned} prove that
unconditional moments of the unpruned second-order solution do not
exist in general, because the quadratic feedback can generate paths
that grow without bound. \citet{denhaan2012solving} further show
that the unpruned dynamics can exhibit ``spurious oscillations'' and
``explosive paths'' that do not correspond to any feature of the true
model solution, but are artifacts of the polynomial approximation
operating outside its radius of convergence. Without stationarity,
there is no well-defined ergodic distribution, and the filtered
state may drift arbitrarily far from the steady state over long
sample periods.

\paragraph{Five structural advantages of pruning for our method.}
The pruned decomposition $x_t = \xf_t + \xs_t$ resolves all of
these issues:

\begin{enumerate}[nosep]
\item \textbf{Guaranteed stationarity of $\xf_t$.}
The first-order state follows linear dynamics
$\xf_{t+1} = \hx\xf_t + \eta\varepsilon_{t+1}$. Under
\Cref{ass:stable} ($\specrad(\hx) < 1$), $\xf_t$ is unconditionally
stationary with well-defined moments. No quadratic term can cause
explosions, regardless of parameter configuration.

\item \textbf{Exact predict step.}
Because the transition of $\xf_t$ is linear, the predict step of
our filter is the \emph{exact} Kalman predict---zero approximation
error. The UKF/CDKF sigma-point approximation is confined to the
observation step, where it handles the \emph{known} quadratic
function $\phi(\xf)$. By contrast, the unpruned UKF/CDKF
introduces sigma-point approximation error in \emph{both} the
predict and update steps, with the predict-step error potentially
compounding across time.

\item \textbf{Exact conditional mean of the observation.}
The observation equation under pruning is \emph{exactly quadratic}
in~$\xf_t$ (\Cref{rem:quad}). The UKF sigma-point mean captures
$\E[\phi(\xf_t) \mid Y_{1:t}]$ with zero error
(\Cref{prop:moment}(a)). Under the unpruned approach, the
observation is quadratic in the \emph{full} state~$x_t$, but $x_t$
includes second-order terms that are correlated with $x_t \kron x_t$
in a non-Gaussian way, degrading the sigma-point accuracy.

\item \textbf{Bounded $\xs_t$ correction.}
The second-order correction satisfies
$\xs_{t+1} = \hx\xs_t + \tfrac{1}{2}\Hxx(\hat\xf_{t|t} \kron \hat\xf_{t|t}) + \tfrac{1}{2}\hsig$.
Since $\specrad(\hx) < 1$ and $\hat\xf_{t|t}$ is bounded (by the
filtered covariance bound of \Cref{prop:stability}), the
correction $\xs_t$ is bounded as well. There is no explosive
feedback.

\item \textbf{Smooth gradients for HMC.}
The log-likelihood under pruning is a smooth function of~$\theta$:
the linear predict step and the polynomial observation step are both
differentiable, and the implicit adjoint gradient through the
perturbation solver is well-defined \citep{childers2024differentiable}.
Under the unpruned approach, the quadratic \emph{transition} creates
gradient instabilities: the derivative of $\Hxx(x \kron x)$ with
respect to~$\theta$ involves products of filtered states with
perturbation-solution derivatives, which can amplify numerical
errors near the stability boundary.
\end{enumerate}

\paragraph{The key conceptual shift.}
Ivashchenko and Andreasen use sigma-point filters to \emph{approximate
the full nonlinear dynamics}---the filter is a substitute for the
intractable optimal nonlinear filter. Our approach is fundamentally
different: we use pruning to \emph{separate the problem} into a
linear filtering component ($\xf_t$) and a deterministic propagation
component ($\xs_t$), and use the UKF only to handle the
\emph{observation nonlinearity}. This separation is what enables
the exact predict step, the exact conditional mean, the stability
guarantee, and the smooth gradients---none of which are available
in the unpruned approach.

\Cref{tab:pruning-comparison} summarizes the comparison.

\begin{table}[htbp]
\centering
\caption{Pruned vs.\ Unpruned Sigma-Point Filtering: Structural Comparison}
\label{tab:pruning-comparison}
\small
\begin{tabular}{@{}p{3.8cm}p{5.2cm}p{5.2cm}@{}}
\toprule
Property
  & Unpruned UKF/CDKF
  & Pruned SR-UKF (this paper) \\
\midrule
Filtered state
  & Full $x_t$ (includes 2nd-order effects)
  & $\xf_t$ only (1st-order) \\
Transition handled by
  & Sigma-point propagation (approximate)
  & Exact Kalman predict (exact) \\
Observation handled by
  & Sigma-point propagation (approximate)
  & Sigma-point propagation (exact mean for quadratic) \\
2nd-order correction
  & Implicit in filtered $x_t$
  & Explicit deterministic $\xs_t$ update \\
Stationarity
  & Not guaranteed (moments may not exist)
  & Guaranteed under $\specrad(\hx) < 1$ \\
Stability under MCMC
  & Can diverge at extreme $\theta$
  & Bounded for all $\theta$ satisfying BK \\
Predict-step error
  & $\cO(\|P\|^{3/2})$ (sigma-point approx.)
  & Zero (exact linear predict) \\
Observation-mean error
  & $\cO(\|P\|^{3/2})$ for quadratic in $x_t$
  & Zero (exact for quadratic in $\xf_t$) \\
Gradient smoothness
  & Unstable near $\specrad(\hx) \to 1$
  & Smooth for all $\specrad(\hx) < 1$ \\
\bottomrule
\end{tabular}
\end{table}

We return to the question of why pruning is essential---and
rebut it in detail for skeptical reviewers---in
\Cref{sec:challenge-unpruned}.

\paragraph{Gap in the literature.}
No existing paper combines the sigma-point methodology of the UKF
with the pruned state-space decomposition of
\citet{andreasen2018pruned}. That is, no prior work applies the
UKF to the first-order state~$\xf_t$ alone while propagating the
second-order correction~$\xs_t$ deterministically, using sigma-point
quadrature to capture the quadratic observation nonlinearity.
\Cref{tab:comparison} summarizes the positioning relative to
all filtering approaches.

\begin{table}[htbp]
\centering
\caption{Comparison of Filtering Approaches for Second-Order DSGE Models}
\label{tab:comparison}
\begin{tabular}{@{}lcccc@{}}
\toprule
Method & State dim.\ & Filter type & Stability & Per-step cost \\
\midrule
Augmented KF
  & $n_x{+}n_x^2{+}n_x$
  & Linear KF
  & Pruning
  & $\cO\!\left((2n_x{+}n_x^2)^3\right)$ \\
UKF (no pruning)
  & $n_x$
  & Sigma-point
  & None
  & $\cO(n_x^3)$ \\
Particle filter
  & $n_x$
  & Monte Carlo
  & N/A
  & $\cO(N_p n_x^2)$ \\
\textbf{Pruned SR-UKF}
  & $\bm{n_x}$
  & \textbf{Sigma-point}
  & \textbf{Pruning}
  & $\bm{\cO(n_x^3)}$ \\
\bottomrule
\end{tabular}
\end{table}

\subsection{Square-Root Sigma-Point Filtering}

The Unscented Transform (UT) of \citet{julier1997new} propagates a
symmetric set of deterministically chosen ``sigma points'' through a
nonlinear function to recover the posterior mean and covariance without
linearization. \citet{wan2000unscented} formalize the Unscented Kalman
Filter (UKF) and demonstrate that it captures the posterior mean and
covariance correctly to third order for Gaussian inputs and arbitrary
nonlinearities---and \emph{exactly} for polynomial nonlinearities of
degree two.

\citet{van2001square} introduce the Square-Root UKF (SR-UKF), which
propagates the Cholesky factor of the covariance matrix directly using
QR decompositions and rank-one Cholesky updates/downdates. The SR-UKF
guarantees positive semi-definiteness of the covariance at every step
(assuming non-negative sigma-point weights), roughly halves the
numerical condition number, and is no more expensive than the standard
UKF.

\subsection{Gradient-Based Bayesian DSGE Estimation}

\citet{childers2024differentiable} develop a framework for
differentiable state-space models, combining implicit
differentiation through the perturbation solver with Hamiltonian
Monte Carlo sampling. Their approach avoids differentiating through
the QZ decomposition and Sylvester equation by instead differentiating
the \emph{implicit equilibrium conditions} that define the
perturbation solution---a technique known as the adjoint or implicit
function theorem method. This provides exact gradients of the
log-posterior with respect to structural parameters at
$\cO(n^2)$ cost per gradient evaluation, enabling the use of HMC
and the No-U-Turn Sampler (NUTS) for Bayesian DSGE estimation.


% ══════════════════════════════════════════════════════════════════════
%  SECTION 3 — MODEL FRAMEWORK
% ══════════════════════════════════════════════════════════════════════
\section{Model Framework}\label{sec:model}

\subsection{General DSGE in Expectations Form}

Consider a DSGE model with $n_y$~control (jump) variables
$y_t \in \R^{n_y}$ and $n_x$~state (predetermined) variables
$x_t \in \R^{n_x}$, subject to $n_\varepsilon$~structural shocks
$\varepsilon_t \sim \mathcal{N}(0, I_{n_\varepsilon})$. The
equilibrium conditions are
\begin{equation}\label{eq:dsge}
\E_t\!\left[
  H\!\left(y_{t+1},\, y_t,\, x_{t+1},\, x_t\right)
\right] = 0,
\end{equation}
where $H : \R^{n_y} \times \R^{n_y} \times \R^{n_x} \times \R^{n_x}
\to \R^{n_y + n_x}$ collects the $n_y + n_x$ equilibrium equations,
and the expectation is conditional on information at time~$t$. The
state transition includes the shock:
$x_{t+1} = \tilde{h}(y_t, x_t) + \eta(\theta)\,\varepsilon_{t+1}$,
where $\eta \in \R^{n_x \times n_\varepsilon}$ is the shock-impact
matrix.

Let $(\bar{y}, \bar{x})$ denote the deterministic steady state
satisfying $H(\bar{y}, \bar{y}, \bar{x}, \bar{x}) = 0$. We work
with deviations from steady state:
$\hat{y}_t \equiv y_t - \bar{y}$, $\hat{x}_t \equiv x_t - \bar{x}$.

\subsection{First-Order Perturbation Solution}\label{sec:first-order}

The first-order (linear) approximation yields policy and transition
functions
\begin{align}
\hat{y}_t &= \gx\,\hat{x}_t, \label{eq:fo-y}\\
\hat{x}_{t+1} &= \hx\,\hat{x}_t + \eta\,\varepsilon_{t+1},
  \label{eq:fo-x}
\end{align}
where $\gx \in \R^{n_y \times n_x}$ and $\hx \in \R^{n_x \times n_x}$
are determined by the first-order conditions. Define the total-variable
vector $v = (y', y, x', x) \in \R^{2(n_y + n_x)}$ and the Jacobians
$H_{y'}, H_y, H_{x'}, H_x$ evaluated at the steady state. The
first-order conditions are:
\begin{equation}\label{eq:fo-implicit}
H_x + H_y\,\gx + (H_{x'} + H_{y'}\,\gx)\,\hx = 0,
\end{equation}
together with the Blanchard--Kahn condition requiring the number of
stable generalized eigenvalues of the pencil $(H_{x'} + H_{y'}, -(H_x + H_y))$
to equal~$n_x$. The matrices $\gx, \hx$ are typically recovered via
the QZ (generalized Schur) decomposition \citep{klein2000using}.

\begin{assumption}\label{ass:bk}
The Blanchard--Kahn condition holds: the generalized Schur
decomposition of the pencil $(\Gamma_1, \Gamma_0)$ with
$\Gamma_0 = [H_{x'}\;|\;H_{y'}]$ and $\Gamma_1 = [-H_x\;|\;-H_y]$
yields exactly $n_x$ stable generalized eigenvalues
(with modulus strictly less than~$1 + \delta$ for some small
$\delta > 0$).
\end{assumption}

\begin{assumption}\label{ass:stable}
The spectral radius of $\hx$ satisfies
$\specrad(\hx) \equiv \max_i |\lambda_i(\hx)| < 1$, ensuring asymptotic
stability of the first-order state dynamics.
\end{assumption}

\subsection{Second-Order Perturbation Solution}\label{sec:second-order}

Following \citet{schmitt2004solving}, the second-order approximation
adds quadratic correction terms:
\begin{align}
\hat{y}_t &= \gx\,\hat{x}_t
  + \tfrac{1}{2}\Gxx\,(\hat{x}_t \kron \hat{x}_t)
  + \tfrac{1}{2}\gss\,\sigma^2,
  \label{eq:so-y} \\
\hat{x}_{t+1} &= \hx\,\hat{x}_t + \eta\,\varepsilon_{t+1}
  + \tfrac{1}{2}\Hxx\,(\hat{x}_t \kron \hat{x}_t)
  + \tfrac{1}{2}\hsig\,\sigma^2,
  \label{eq:so-x}
\end{align}
where $\Gxx \in \R^{n_y \times n_x^2}$,
$\Hxx \in \R^{n_x \times n_x^2}$,
$\gss \in \R^{n_y}$, and $\hsig \in \R^{n_x}$ are the second-order
coefficient matrices. Throughout, we set $\sigma = 1$ (full
stochasticity) so that $\gss$ and $\hsig$ are vectors.

\paragraph{The Sylvester equation for $\Gxx, \Hxx$.}
Define $n \equiv n_y + n_x$,
$M \equiv [\gx\hx;\;\gx;\;\hx;\;I_{n_x}] \in \R^{2n \times n_x}$,
and the Hessian tensor
$\Psi \in \R^{n \times 2n \times 2n}$ where
$\Psi_{ijk} = \partial^2 H_i / (\partial v_j\,\partial v_k)$
evaluated at steady state. The stacked second-order coefficient
$X \equiv [\Gxx;\;\Hxx] \in \R^{n \times n_x^2}$ solves the
generalized Sylvester equation
\begin{equation}\label{eq:sylvester}
A\,X\,(\hx \kron \hx)^\top + D\,X = -C,
\end{equation}
where
$A = [H_{y'}\;\;0_{n_y \times n_x};\;0_{n_x \times n_y}\;\;0_{n_x}]$,
$D = [H_y\;\;H_{y'}\gx + H_{x'}]$,
and $C \in \R^{n \times n_x^2}$ is formed from the quadratic terms
$C_{i \cdot} = \vect(M^\top \Psi_i M)^\top$ for each equation~$i$.
The Sylvester equation~\eqref{eq:sylvester} is solved efficiently
by the Bartels--Stewart algorithm: Schur decomposition of~$\hx$ gives
$\hx = Q T Q^\top$, which induces a block-triangular structure in
$\hx \kron \hx$, reducing the problem to a sequence of
$n \times n$ linear systems \citep{bartels1972solution}.

\paragraph{The volatility corrections $\gss, \hsig$.}
The constants $\gss, \hsig$ solve a linear system
\begin{equation}\label{eq:volatility}
\begin{bmatrix} H_y & H_{y'}\gx + H_{x'} \end{bmatrix}
\begin{bmatrix} \gss \\ \hsig \end{bmatrix}
= -B,
\end{equation}
where $B \in \R^n$ depends on $\eta$, the Hessians~$\Psi$, and the
second-order coefficients $\Gxx, \Hxx$ already obtained.

\subsection{The Pruned State-Space System}\label{sec:pruned-ss}

Direct simulation of~\eqref{eq:so-x} feeds
$\hat{x}_t \kron \hat{x}_t$ (including second-order terms) back into
the quadratic dynamics, creating a feedback loop that can generate
explosive paths even when $\specrad(\hx) < 1$.
\citet{kim2008calculating} and \citet{andreasen2018pruned} resolve
this by decomposing the state:

\begin{definition}[Pruned state-space system]\label{def:pruned}
Define the first-order state~$\xf_t$ and the second-order
correction~$\xs_t$ by:
\begin{align}
\xf_{t+1} &= \hx\,\xf_t + \eta\,\varepsilon_{t+1},
  \label{eq:pruned-xf} \\
\xs_{t+1} &= \hx\,\xs_t
  + \tfrac{1}{2}\Hxx\,(\xf_t \kron \xf_t)
  + \tfrac{1}{2}\hsig,
  \label{eq:pruned-xs} \\
y_t &= \gx\,(\xf_t + \xs_t)
  + \tfrac{1}{2}\Gxx\,(\xf_t \kron \xf_t)
  + \tfrac{1}{2}\gss + \varepsilon_t^y,
  \label{eq:pruned-obs}
\end{align}
where $\varepsilon_t^y \sim \mathcal{N}(0, R)$ is measurement noise.
\end{definition}

\noindent
The key feature of pruning is that \emph{only the first-order
state~$\xf_t$ enters the quadratic terms}. Since $\xf_t$ follows
linear-Gaussian dynamics~\eqref{eq:pruned-xf}, the Kronecker product
$\xf_t \kron \xf_t$ remains bounded in expectation, preventing
explosive behavior.

\subsection{The Augmented State-Space (Existing Approach)}\label{sec:augmented}

\citet{andreasen2018pruned} and \citet{kollmann2014tractable} define
the augmented state
\begin{equation}\label{eq:aug-state}
\xi_t \equiv
\begin{bmatrix} \xf_t \\ \vect(\xf_t {\xf_t}^\top) \\ \xs_t
\end{bmatrix}
\in \R^{n_x + n_x^2 + n_x}
\end{equation}
and show that the system~\eqref{eq:pruned-xf}--\eqref{eq:pruned-obs}
can be written as a linear state-space model:
\begin{align}
\xi_{t+1} &= \mathcal{T}\,\xi_t + c + \nu_{t+1},
  \quad \E[\nu_t\nu_t^\top] = \mathcal{Q},
  \label{eq:aug-transition}\\
y_t &= \mathcal{Z}\,\xi_t + d + \varepsilon_t^y,
  \label{eq:aug-obs}
\end{align}
where
\begin{equation}\label{eq:aug-T}
\mathcal{T} =
\begin{bmatrix}
\hx & 0 & 0 \\
0 & \hx \kron \hx & 0 \\
0 & \frac{1}{2}\Hxx & \hx
\end{bmatrix},
\quad
c =
\begin{bmatrix}
0 \\
\vect(\eta\eta^\top) \\
\frac{1}{2}\hsig
\end{bmatrix},
\end{equation}
and $\mathcal{Z}$ and~$d$ select and combine the relevant rows to
form the observation equation. The process-noise covariance
$\mathcal{Q}$ is block-diagonal with the $(1,1)$~block equal to
$\eta\eta^\top$ and the $(2,2)$~block $\mathcal{Q}_{kk}$ given by
the \emph{unconditional} covariance of the Kronecker innovations
\citep[eq.~16]{andreasen2018pruned}:
\begin{equation}\label{eq:Q-kk}
\mathcal{Q}_{kk} =
(\Sigma_a \kron \Sigma_b)
+ (\Sigma_b \kron \Sigma_a)
+ (P_1 \kron P_1)
+ \text{cross terms},
\end{equation}
where $\Sigma_a = \hx P_f \hx^\top$, $\Sigma_b = \eta\eta^\top$,
$P_1 = \eta\eta^\top$, and $P_f$ is the unconditional covariance of
$\xf_t$ solving the discrete Lyapunov equation
$P_f = \hx P_f \hx^\top + \eta\eta^\top$.

The standard Kalman filter applied to the augmented
system~\eqref{eq:aug-transition}--\eqref{eq:aug-obs} yields an
exact (under the Gaussian approximation) log-likelihood, but the
per-step cost is
\begin{equation}\label{eq:aug-cost}
\cO\!\left((n_x + n_x^2 + n_x)^3\right) = \cO(n_x^6),
\end{equation}
since the dominant operations (Kalman gain, covariance update) involve
matrices of dimension $2n_x + n_x^2$.

\begin{remark}
For the \citet{schmitt2003closing} model with $n_x = 7$, the
augmented dimension is $7 + 49 + 7 = 63$. The Kalman filter operates
on $63 \times 63$ matrices---manageable but expensive when embedded
in an MCMC sampler requiring $10^4$--$10^5$ likelihood evaluations.
For a medium-scale New Keynesian model with $n_x = 30$, the augmented
dimension is~960, making the approach infeasible.
\end{remark}


% ══════════════════════════════════════════════════════════════════════
%  SECTION 4 — THE PRUNED SR-UKF
% ══════════════════════════════════════════════════════════════════════
\section{The Pruned Square-Root Unscented Kalman Filter}\label{sec:method}

\subsection{Overview and Intuition}

The central idea of the Pruned SR-UKF is to exploit the structure of
the pruned system~\eqref{eq:pruned-xf}--\eqref{eq:pruned-obs}:
\begin{enumerate}[nosep]
\item The \emph{transition} of $\xf_t$ is \emph{linear}:
$\xf_{t+1} = \hx\,\xf_t + \eta\,\varepsilon_{t+1}$.
No sigma-point propagation is needed for the predict step---the
Kalman filter predict equations apply exactly.
\item The \emph{observation} equation is \emph{quadratic} in~$\xf_t$
(with $\xs_t$ entering linearly as a known offset). The unscented
transform handles this nonlinearity by propagating $2n_x + 1$
sigma points through the observation function.
\item The \emph{second-order correction}~$\xs_t$ depends on the
Kronecker product $\xf_t \kron \xf_t$, which under Gaussian
filtering has conditional expectation
$\E[\xf_t \kron \xf_t \mid Y_{1:t}]
= \hat{\xf}_{t|t} \kron \hat{\xf}_{t|t} + \vect(P_{t|t})$.
We approximate this by $\hat{\xf}_{t|t} \kron \hat{\xf}_{t|t}$
alone---the ``certainty equivalence'' approximation---incurring
an $\cO(\sigma^2)$ error that we characterize in
\Cref{sec:theory}.
\end{enumerate}
The filter state is $(\hat{\xf}_{t|t}, S_{t|t}, \xs_t)$, where
$S_{t|t}$ is the Cholesky factor of the filtered covariance
$P_{t|t} = S_{t|t} S_{t|t}^\top$. The dimension of the covariance
is $n_x \times n_x$, not $(2n_x + n_x^2) \times (2n_x + n_x^2)$.

\subsection{Sigma-Point Construction}\label{sec:sigma}

Let $L \equiv n_x$ be the state dimension. Choose tuning parameters
$\alpha > 0$, $\beta \geq 0$, $\kappa \geq 0$, and define
\begin{equation}\label{eq:lambda}
\lambda = \alpha^2(L + \kappa) - L, \qquad
\gamma = \sqrt{L + \lambda}.
\end{equation}
Given the state estimate $\hat{x} \in \R^L$ and its Cholesky factor
$S \in \R^{L \times L}$ (so $P = SS^\top$), define $2L + 1$
sigma points:
\begin{align}
\chi^{(0)} &= \hat{x}, \label{eq:sp0}\\
\chi^{(i)} &= \hat{x} + \gamma\,S_{:,i},
  \quad i = 1, \ldots, L, \label{eq:sp-plus}\\
\chi^{(L+i)} &= \hat{x} - \gamma\,S_{:,i},
  \quad i = 1, \ldots, L, \label{eq:sp-minus}
\end{align}
where $S_{:,i}$ denotes the $i$-th column of~$S$. The associated
weights are:
\begin{align}
w_m^{(0)} &= \frac{\lambda}{L + \lambda}, &
w_c^{(0)} &= w_m^{(0)} + (1 - \alpha^2 + \beta),
  \label{eq:w0} \\
w_m^{(j)} &= w_c^{(j)} = \frac{1}{2(L + \lambda)},
  &j &= 1, \ldots, 2L. \label{eq:wj}
\end{align}

We use the default parameters $\alpha = 1$, $\beta = 2$,
$\kappa = \max(0, 3 - L)$, which ensure $w_c^{(0)} \geq 0$ for
all~$L \geq 1$ and provide exact recovery of the Gaussian mean and
covariance for affine transformations.

\subsection{Initialization}\label{sec:init}

The filter is initialized at the unconditional moments of the
first-order state:
\begin{equation}\label{eq:init}
\hat{\xf}_{0|0} = 0, \qquad
P_{0|0} = P_f + \epsilon\,I_{n_x}, \qquad
\xs_0 = 0,
\end{equation}
where $P_f$ solves the discrete Lyapunov equation
\begin{equation}\label{eq:lyapunov}
P_f = \hx\,P_f\,\hx^\top + \eta\eta^\top
\end{equation}
and $\epsilon > 0$ is a small nugget (e.g., $10^{-10}$) ensuring
positive definiteness.

The Lyapunov equation is solved iteratively: starting from
$P^{(0)} = \eta\eta^\top$, iterate
$P^{(k+1)} = \hx P^{(k)} \hx^\top + \eta\eta^\top$ until
$\|P^{(k+1)} - P^{(k)}\|_F / \|P^{(k+1)}\|_F < \tau$
with tolerance $\tau = 10^{-12}$.

\subsection{Predict Step}\label{sec:predict}

Since the transition $\xf_{t+1} = \hx\,\xf_t + \eta\,\varepsilon_{t+1}$
is linear, the predict step is exact:
\begin{align}
\hat{\xf}_{t+1|t} &= \hx\,\hat{\xf}_{t|t},
  \label{eq:predict-mean} \\
P_{t+1|t} &= \hx\,P_{t|t}\,\hx^\top + Q_f,
  \label{eq:predict-cov}
\end{align}
where $Q_f = \eta\eta^\top + \epsilon\,I_{n_x}$.

In square-root form, we compute $S_{t+1|t}$ from the QR decomposition
of the matrix
\begin{equation}\label{eq:qr-predict}
\mathcal{R}_p =
\begin{bmatrix}
\sqrt{w_c^{(0)}}\,(\hx\,S_{t|t})^\top \\
\vdots \\
S_{Q_f}^\top
\end{bmatrix},
\end{equation}
where $S_{Q_f} = \chol(Q_f)$, so that $S_{t+1|t}$ is the $R$ factor
of the QR decomposition (appropriately transposed and sign-corrected).

\begin{remark}
Because the transition is linear, the sigma-point propagation
$\chi^{(j)}_+ = \hx\,\chi^{(j)}$ and the weighted-mean formula
$\sum_j w_m^{(j)} \chi^{(j)}_+ = \hx\,\hat{x}$ are exactly the
Kalman predict. The QR form is used for numerical stability, not
for handling nonlinearity.
\end{remark}

\subsection{Observation Update}\label{sec:update}

\subsubsection{Sigma-point propagation through the observation function.}

Generate fresh sigma points $\{\tilde{\chi}^{(j)}\}_{j=0}^{2L}$ from
$(\hat{\xf}_{t+1|t},\;S_{t+1|t})$ using
\eqref{eq:sp0}--\eqref{eq:sp-minus}. Define the observation function
\begin{equation}\label{eq:obs-fn}
\phi(x) \equiv
  g_{\text{obs}}\,(x + \xs_t)
  + \tfrac{1}{2}\,G_{\text{obs}}\,(x \kron x)
  + \tfrac{1}{2}\,g_{\text{obs},\sigma\sigma},
\end{equation}
where $g_{\text{obs}} \in \R^{n_y \times n_x}$,
$G_{\text{obs}} \in \R^{n_y \times n_x^2}$, and
$g_{\text{obs},\sigma\sigma} \in \R^{n_y}$ are the observation-mapped
rows of $\gx$, $\Gxx$, and $\gss$ (selecting observed controls and
states). Propagate:
\begin{equation}\label{eq:obs-prop}
z^{(j)} = \phi(\tilde{\chi}^{(j)}),
\quad j = 0, 1, \ldots, 2L.
\end{equation}

\begin{remark}\label{rem:quad}
The function $\phi$ is a polynomial of degree two in its
argument~$x$: it is affine in~$x$ (through $g_{\text{obs}}(x + \xs_t)$)
plus quadratic (through $G_{\text{obs}}(x \kron x)$). The second-order
correction~$\xs_t$ enters only linearly, as a known shift.
\end{remark}

\subsubsection{Innovation statistics.}

Predicted measurement mean:
\begin{equation}\label{eq:z-bar}
\bar{z}_{t+1} = \sum_{j=0}^{2L} w_m^{(j)}\,z^{(j)}.
\end{equation}
Innovation:
\begin{equation}
v_{t+1} = y_{t+1} - \bar{z}_{t+1}.
\end{equation}
Innovation covariance (square-root form): compute $S_{yy}$ from the
QR decomposition of
\begin{equation}\label{eq:qr-obs}
\mathcal{R}_y =
\begin{bmatrix}
\sqrt{w_c^{(0)}}\,(z^{(0)} - \bar{z}_{t+1})^\top \\
\sqrt{w_c^{(1)}}\,(z^{(1)} - \bar{z}_{t+1})^\top \\
\vdots \\
\sqrt{w_c^{(2L)}}\,(z^{(2L)} - \bar{z}_{t+1})^\top \\
S_R^\top
\end{bmatrix},
\end{equation}
where $S_R = \chol(R)$ is the Cholesky factor of the measurement noise.
Then $P_{yy} = S_{yy} S_{yy}^\top$.

\subsubsection{Cross-covariance and Kalman gain.}

The state--measurement cross-covariance is
\begin{equation}\label{eq:Pxz}
P_{xz} = \sum_{j=0}^{2L} w_c^{(j)}\,
  (\tilde{\chi}^{(j)} - \hat{\xf}_{t+1|t})\,
  (z^{(j)} - \bar{z}_{t+1})^\top.
\end{equation}
The Kalman gain is
\begin{equation}\label{eq:gain}
K = P_{xz}\,P_{yy}^{-1}
  = P_{xz}\,S_{yy}^{-\top}\,S_{yy}^{-1},
\end{equation}
computed via two triangular solves.

\subsubsection{Filtered state and covariance.}

The filtered mean is
\begin{equation}\label{eq:filter-mean}
\hat{\xf}_{t+1|t+1} = \hat{\xf}_{t+1|t} + K\,v_{t+1}.
\end{equation}
The filtered covariance is updated as
\begin{equation}\label{eq:filter-cov}
P_{t+1|t+1} = P_{t+1|t} - K\,P_{yy}\,K^\top + \epsilon\,I_{n_x},
\end{equation}
symmetrized via $P \leftarrow \tfrac{1}{2}(P + P^\top)$, and the
Cholesky factor $S_{t+1|t+1} = \chol(P_{t+1|t+1})$ is recomputed.
The nugget $\epsilon > 0$ ensures positive definiteness.

\begin{remark}[On the covariance update form]
The classical SR-UKF of \citet{van2001square} uses rank-one Cholesky
downdates to maintain the square-root factor directly: $S_u$ is
obtained by successively downdating $S_p$ with the columns of
$K\,S_{yy}$. While this guarantees PSD preservation under exact
arithmetic and non-negative weights, the rank-one downdate can fail
in finite precision when $K P_{yy} K^\top$ approaches $P_{t+1|t}$
(near-degenerate innovation). We adopt the computationally equivalent
approach of~\eqref{eq:filter-cov} with nugget regularization, which
is robust in automatic-differentiation frameworks where Cholesky
downdates are not readily available. See \Cref{sec:challenge-sr} for
a detailed discussion.
\end{remark}

\subsection{Deterministic $\xs$ Propagation}\label{sec:xs-update}

After computing the filtered mean $\hat{\xf}_{t+1|t+1}$, the
second-order correction is updated deterministically:
\begin{equation}\label{eq:xs-update}
\xs_{t+1} = \hx\,\xs_t
  + \tfrac{1}{2}\Hxx\,
    (\hat{\xf}_{t+1|t+1} \kron \hat{\xf}_{t+1|t+1})
  + \tfrac{1}{2}\hsig.
\end{equation}

This equation uses the \emph{filtered} (posterior) mean of~$\xf_{t+1}$
in the Kronecker product, rather than the conditional expectation
$\E[\xf_{t+1} \kron \xf_{t+1} \mid Y_{1:t+1}]$. The latter equals
$\hat{\xf}_{t+1|t+1} \kron \hat{\xf}_{t+1|t+1}
+ \vect(P_{t+1|t+1})$;
the omitted term $\tfrac{1}{2}\Hxx\,\vect(P_{t+1|t+1})$ constitutes
the approximation bias, which we analyze in \Cref{sec:bias}.

\subsection{Log-Likelihood}\label{sec:loglik}

The Gaussian log-likelihood contribution at time $t+1$ is
\begin{equation}\label{eq:loglik}
\ell_{t+1} = -\frac{n_y}{2}\ln(2\pi)
  - \ln|\det(S_{yy})|
  - \frac{1}{2}\|S_{yy}^{-1}\,v_{t+1}\|^2,
\end{equation}
where $\ln|\det(S_{yy})| = \sum_{i=1}^{n_y} \ln|S_{yy,ii}|$. The
total log-likelihood is
$\ell(Y_{1:T} \mid \theta) = \sum_{t=1}^{T} \ell_t$.

\subsection{Complete Algorithm}

\Cref{alg:pruned-sr-ukf} summarizes the Pruned SR-UKF.

\begin{algorithm}[htbp]
\caption{Pruned Square-Root UKF for Second-Order DSGE}
\label{alg:pruned-sr-ukf}
\begin{algorithmic}[1]
\Require Model solution
$(\gx, \hx, \eta, \Gxx, \Hxx, \gss, \hsig)$;
observation matrices $(g_{\text{obs}}, G_{\text{obs}}, g_{\text{obs},\sigma\sigma})$;
data $Y = \{y_1, \ldots, y_T\}$;
measurement noise $R$.
\Ensure Log-likelihood $\ell$.
\State Solve Lyapunov equation for $P_f$; set
$S_0 \leftarrow \chol(P_f + \epsilon I)$,
$\hat{\xf}_0 \leftarrow 0$,
$\xs_0 \leftarrow 0$,
$\ell \leftarrow 0$.
\For{$t = 0, 1, \ldots, T-1$}
  \State \textbf{Predict:}
    $\hat{\xf}^- \leftarrow \hx\,\hat{\xf}_t$; \quad
    $S^- \leftarrow \text{QR-predict}(\hx, S_t, S_{Q_f})$.
  \State \textbf{Sigma points:}
    Generate $\{\tilde\chi^{(j)}\}$ from $(\hat{\xf}^-, S^-)$.
  \State \textbf{Observation propagation:}
    $z^{(j)} \leftarrow \phi(\tilde\chi^{(j)})$ for each~$j$.
  \State \textbf{Innovation:}
    $\bar{z} \leftarrow \sum w_m^{(j)} z^{(j)}$; \quad
    $v \leftarrow y_{t+1} - \bar{z}$.
  \State \textbf{Innovation covariance:}
    $S_{yy} \leftarrow \text{QR-obs}(\{z^{(j)}\}, \bar{z}, S_R)$.
  \State \textbf{Cross-covariance:}
    $P_{xz} \leftarrow \sum w_c^{(j)}
      (\tilde\chi^{(j)} - \hat{\xf}^-)
      (z^{(j)} - \bar{z})^\top$.
  \State \textbf{Kalman gain:}
    $K \leftarrow P_{xz}\,(S_{yy} S_{yy}^\top)^{-1}$.
  \State \textbf{Filter update:}
    $\hat{\xf}_{t+1} \leftarrow \hat{\xf}^- + K v$; \quad
    $P^+ \leftarrow S^- {S^-}^\top - K S_{yy} {S_{yy}}^\top K^\top
      + \epsilon I$; \quad
    $S_{t+1} \leftarrow \chol(\tfrac{1}{2}(P^+ + {P^+}^\top))$.
  \State \textbf{$\xs$ update:}
    $\xs_{t+1} \leftarrow \hx\,\xs_t
      + \frac{1}{2}\Hxx(\hat{\xf}_{t+1} \kron \hat{\xf}_{t+1})
      + \frac{1}{2}\hsig$.
  \State \textbf{Log-likelihood:}
    $\ell \leftarrow \ell - \frac{n_y}{2}\ln(2\pi)
      - \sum_i \ln|S_{yy,ii}|
      - \frac{1}{2}\|S_{yy}^{-1} v\|^2$.
\EndFor
\State \Return $\ell$.
\end{algorithmic}
\end{algorithm}


% ══════════════════════════════════════════════════════════════════════
%  SECTION 5 — THEORETICAL PROPERTIES
% ══════════════════════════════════════════════════════════════════════
\section{Theoretical Properties}\label{sec:theory}

In this section we establish the key analytical properties of the
Pruned SR-UKF: moment matching of the observation equation
(\Cref{prop:moment}), characterization of the $\xs$ approximation
bias (\Cref{prop:bias}), computational complexity (\Cref{prop:cost}),
stability (\Cref{prop:stability}), and two limiting-case theorems
(\Cref{thm:linear,thm:zero-cov}).

\subsection{Moment Matching for the Quadratic Observation}

\begin{proposition}[Exact mean, second-order covariance]
\label{prop:moment}
Let $\xf \sim \mathcal{N}(\hat{x}, P)$ with
$P = SS^\top$, and define the quadratic observation function
\begin{equation}\label{eq:phi}
\phi(x) = g_{\mathrm{obs}}(x + \xs)
  + \tfrac{1}{2}\,G_{\mathrm{obs}}\,(x \kron x)
  + \tfrac{1}{2}\,g_{\mathrm{obs},\sigma\sigma},
\end{equation}
where $\xs \in \R^{n_x}$ is a fixed offset. Then the sigma-point
approximations satisfy:
\begin{enumerate}[label=(\alph*)]
\item The predicted measurement mean
$\bar{z} = \sum_{j=0}^{2L} w_m^{(j)}\,\phi(\chi^{(j)})$
equals $\E[\phi(\xf)]$ \emph{exactly}:
\begin{equation}\label{eq:exact-mean}
\bar{z} = g_{\mathrm{obs}}(\hat{x} + \xs)
  + \tfrac{1}{2}\,G_{\mathrm{obs}}\,
    (\hat{x} \kron \hat{x} + \vect(P))
  + \tfrac{1}{2}\,g_{\mathrm{obs},\sigma\sigma}.
\end{equation}
\item The predicted measurement covariance
$P_{zz} = \sum_{j=0}^{2L} w_c^{(j)}\,(z^{(j)} - \bar{z})
(z^{(j)} - \bar{z})^\top$
agrees with $\Var[\phi(\xf)]$ to second order in~$P$:
\begin{equation}\label{eq:cov-approx}
P_{zz} = \Var[\phi(\xf)] + \cO(\|P\|^2).
\end{equation}
\end{enumerate}
\end{proposition}

\begin{proof}
\textbf{Part (a).}
Write $\phi(x) = a + B\,x + \tfrac{1}{2}\,G_{\mathrm{obs}}\,(x \kron x)$,
where $a = g_{\mathrm{obs}}\,\xs + \tfrac{1}{2}\,g_{\mathrm{obs},\sigma\sigma}$
and $B = g_{\mathrm{obs}}$.

The expectation is
\begin{equation}
\E[\phi(\xf)] = a + B\,\hat{x}
  + \tfrac{1}{2}\,G_{\mathrm{obs}}\,
    \E[\xf \kron \xf].
\end{equation}
For $\xf \sim \mathcal{N}(\hat{x}, P)$, we have
$\E[\xf \kron \xf] = \hat{x} \kron \hat{x} + \vect(P)$, giving
\eqref{eq:exact-mean}.

Now consider the sigma-point mean. By the symmetry of the sigma
points, $\sum_j w_m^{(j)} \chi^{(j)} = \hat{x}$ (the UT preserves
the mean for any function). For the quadratic term, using
$\chi^{(i)} = \hat{x} + \gamma\,S_{:,i}$ and
$\chi^{(L+i)} = \hat{x} - \gamma\,S_{:,i}$:
\begin{align}
\sum_{j=0}^{2L} &w_m^{(j)}\,(\chi^{(j)} \kron \chi^{(j)})
\notag\\
&= w_m^{(0)}\,(\hat{x} \kron \hat{x})
+ \sum_{i=1}^{L} w_m^{(i)}\bigl[
  (\hat{x} + \gamma S_{:,i}) \kron (\hat{x} + \gamma S_{:,i})
  + (\hat{x} - \gamma S_{:,i}) \kron (\hat{x} - \gamma S_{:,i})
\bigr] \notag\\
&= w_m^{(0)}\,(\hat{x} \kron \hat{x})
+ \sum_{i=1}^{L} w_m^{(i)}\bigl[
  2(\hat{x} \kron \hat{x})
  + 2\gamma^2\,(S_{:,i} \kron S_{:,i})
\bigr] \notag\\
&= \Bigl(w_m^{(0)} + 2L\,w_m^{(1)}\Bigr)(\hat{x} \kron \hat{x})
+ 2\gamma^2 w_m^{(1)} \sum_{i=1}^{L} (S_{:,i} \kron S_{:,i}).
\label{eq:kron-sigma}
\end{align}
From \eqref{eq:w0}--\eqref{eq:wj}:
$w_m^{(0)} + 2L\,w_m^{(1)}
= \frac{\lambda}{L+\lambda} + \frac{2L}{2(L+\lambda)} = 1$.
Also, $2\gamma^2 w_m^{(1)}
= 2(L+\lambda) \cdot \frac{1}{2(L+\lambda)} = 1$,
and $\sum_{i=1}^{L} S_{:,i} \kron S_{:,i} = \vect(SS^\top) = \vect(P)$.
Therefore
\begin{equation}
\sum_{j=0}^{2L} w_m^{(j)}\,(\chi^{(j)} \kron \chi^{(j)})
= \hat{x} \kron \hat{x} + \vect(P),
\end{equation}
and $\bar{z} = a + B\hat{x}
+ \tfrac{1}{2}G_{\mathrm{obs}}(\hat{x} \kron \hat{x} + \vect(P))
= \E[\phi(\xf)]$. \hfill$\square$ (Part a)

\medskip
\textbf{Part (b).}
The true variance of $\phi(\xf) = a + B\xf + \tfrac{1}{2}G_{\mathrm{obs}}(\xf \kron \xf)$
involves moments of $\xf$ up to fourth order. For Gaussian $\xf$,
the fourth central moments are determined by $P$ via
$\E[(\Delta x)^{\kron 4}]
= \vect(P) \vect(P)^\top + (I + K_{n_x,n_x})(P \kron P)$,
where $K_{n_x,n_x}$ is the commutation matrix.

The sigma-point covariance captures the leading terms
(those linear in~$P$) exactly. The error arises from the fourth-order
terms, which contribute $\cO(\|P\|^2)$. Specifically, the UKF
covariance approximation for a degree-$d$ polynomial is exact through
order $2d - 1$ in the standard deviations when $\alpha = 1$ and
$\kappa = 3 - L$ \citep{julier2004unscented}. For $d = 2$, this
means exactness through third order, implying the error in the
covariance is fourth order in the standard deviations, or equivalently
$\cO(\|P\|^2)$.
\end{proof}

\subsection{Bias Characterization of the $\xs$ Update}\label{sec:bias}

\begin{proposition}[$\xs$ approximation bias]\label{prop:bias}
Let $\hat{\xf}_{t|t}$ and $P_{t|t}$ be the filtered mean and
covariance of~$\xf_t$ given observations $Y_{1:t}$, and let
$\xs_t^{\star}$ denote the Bayesian optimal second-order correction
obtained by using
$\E[\xf_t \kron \xf_t \mid Y_{1:t}]
= \hat{\xf}_{t|t} \kron \hat{\xf}_{t|t} + \vect(P_{t|t})$
in~\eqref{eq:pruned-xs}. The Pruned SR-UKF
correction~$\xs_{t+1}^{\mathrm{UKF}}$
from~\eqref{eq:xs-update} satisfies
\begin{equation}\label{eq:xs-bias}
\xs_{t+1}^{\star} - \xs_{t+1}^{\mathrm{UKF}}
= \tfrac{1}{2}\,\Hxx\,\vect(P_{t|t})
+ \hx\,(\xs_t^{\star} - \xs_t^{\mathrm{UKF}}).
\end{equation}
In particular, the per-step bias increment is
$\Delta_t \equiv \tfrac{1}{2}\Hxx\,\vect(P_{t|t})$, and the
cumulative bias follows
\begin{equation}\label{eq:cum-bias}
\xs_t^{\star} - \xs_t^{\mathrm{UKF}}
= \sum_{s=0}^{t-1} \hx^{t-1-s}\,\Delta_s.
\end{equation}
\end{proposition}

\begin{proof}
The Bayesian optimal update is
\begin{equation}
\xs_{t+1}^{\star} = \hx\,\xs_t^{\star}
  + \tfrac{1}{2}\Hxx\,
    (\hat{\xf}_{t|t} \kron \hat{\xf}_{t|t} + \vect(P_{t|t}))
  + \tfrac{1}{2}\hsig.
\end{equation}
The Pruned SR-UKF update is
\begin{equation}
\xs_{t+1}^{\mathrm{UKF}} = \hx\,\xs_t^{\mathrm{UKF}}
  + \tfrac{1}{2}\Hxx\,
    (\hat{\xf}_{t|t} \kron \hat{\xf}_{t|t})
  + \tfrac{1}{2}\hsig.
\end{equation}
Subtracting yields~\eqref{eq:xs-bias}. Unrolling the recursion with
$\xs_0^{\star} = \xs_0^{\mathrm{UKF}} = 0$
gives~\eqref{eq:cum-bias}.
\end{proof}

\begin{corollary}\label{cor:bias-bound}
Under \Cref{ass:stable}, the cumulative bias is bounded:
\begin{equation}\label{eq:bias-bound}
\bigl\|\xs_t^{\star} - \xs_t^{\mathrm{UKF}}\bigr\|
\leq \frac{\|\Hxx\|}{2(1 - \specrad(\hx))}\,
  \sup_{0 \leq s < t} \bigl\|\vect(P_{s|s})\bigr\|.
\end{equation}
Since $P_{s|s}$ is bounded and $\|P_{s|s}\| = \cO(\sigma^2)$
(where $\sigma$ scales the shock standard deviations), the
cumulative bias is $\cO(\sigma^2)$.
\end{corollary}

\begin{proof}
From \eqref{eq:cum-bias},
$\|\xs_t^{\star} - \xs_t^{\mathrm{UKF}}\|
\leq \sum_{s=0}^{t-1} \|\hx\|^{t-1-s} \|\Delta_s\|
\leq \frac{\|\Hxx\|}{2}\,\sup_s \|\vect(P_{s|s})\|
  \,\sum_{k=0}^{\infty} \specrad(\hx)^k
= \frac{\|\Hxx\|}{2(1-\specrad(\hx))}\,\sup_s \|\vect(P_{s|s})\|$.
\end{proof}

\begin{remark}[Numerical magnitude]
For the \citet{schmitt2003closing} model with shock standard deviation
$\sigma_e = 0.0129$, we have $\|P_{t|t}\| \approx 10^{-4}$ and
$\|\Hxx\| \approx \cO(1)$, giving a per-step bias of approximately
$5 \times 10^{-5}$. This is far below the measurement noise level
and the Monte Carlo error of any posterior sampler.
\end{remark}

\subsection{Computational Complexity}

\begin{proposition}[Per-step cost]\label{prop:cost}
The per-step cost of the Pruned SR-UKF is
$\cO(n_x^2 n_y + n_x^3)$, while the augmented Kalman filter requires
$\cO((2n_x + n_x^2)^3)$ operations. The speedup factor is
$\cO(n_x^3)$.
\end{proposition}

\begin{proof}
\textbf{Pruned SR-UKF.}
The dominant operations per step are:
\begin{itemize}[nosep]
\item Sigma-point generation: $\cO(n_x^2)$ (scaling columns of~$S$).
\item Sigma-point propagation through $\phi$: $2n_x + 1$ evaluations,
each requiring a matrix-vector product $g_{\mathrm{obs}} x$
($\cO(n_y n_x)$) and the Kronecker product $G_{\mathrm{obs}}(x \kron x)$
($\cO(n_y n_x^2)$). Total: $\cO(n_x \cdot n_y n_x^2) = \cO(n_y n_x^3)$.
\item QR factorization for $S_{yy}$: the matrix has
$(2n_x + 1 + n_y)$ rows and $n_y$ columns;
cost is $\cO(n_x n_y^2)$.
\item Cross-covariance $P_{xz}$: $\cO(n_x^2 n_y)$.
\item Kalman gain via triangular solves: $\cO(n_x n_y^2)$.
\item Covariance update and Cholesky: $\cO(n_x^3)$.
\item $\xs$ update: $\cO(n_x^3)$ (for $\Hxx(x \kron x)$).
\end{itemize}
The leading term is $\cO(n_y n_x^3)$; since typically $n_y \leq n_x$,
this simplifies to $\cO(n_x^3)$ (with a constant depending on~$n_y$).

\textbf{Augmented KF.}
The augmented state dimension is $n_a = 2n_x + n_x^2$. The Kalman
filter requires:
\begin{itemize}[nosep]
\item Predicted covariance: $\mathcal{T} P \mathcal{T}^\top + \mathcal{Q}$,
which is $\cO(n_a^2 \cdot n_a) = \cO(n_a^3)$.
\item Kalman gain: $\cO(n_a^2 n_y + n_y^3)$.
\item Covariance update: $\cO(n_a^2 n_y)$.
\end{itemize}
Total: $\cO(n_a^3) = \cO((2n_x + n_x^2)^3)$. For large~$n_x$,
$n_a \approx n_x^2$ and the cost is $\cO(n_x^6)$.

The ratio is $\cO(n_x^6 / n_x^3) = \cO(n_x^3)$.
\end{proof}

\begin{table}[htbp]
\centering
\caption{Computational Cost Comparison (Operations per Time Step)}
\label{tab:cost}
\begin{tabular}{@{}lrrr@{}}
\toprule
$n_x$ & Aug.\ dim.\ $n_a$ & Aug.\ KF cost
  & Pruned SR-UKF cost \\
\midrule
 7 &  63 & $2.5 \times 10^5$ & $3.4 \times 10^2$ \\
15 & 255 & $1.7 \times 10^7$ & $3.4 \times 10^3$ \\
30 & 960 & $8.8 \times 10^8$ & $2.7 \times 10^4$ \\
50 & 2600 & $1.8 \times 10^{10}$ & $1.3 \times 10^5$ \\
\bottomrule
\end{tabular}
\end{table}

\subsection{Stability}\label{sec:stability}

\begin{proposition}[Covariance boundedness]\label{prop:stability}
Under \Cref{ass:stable} and the conditions:
\begin{enumerate}[label=(\roman*),nosep]
\item $(h_x, \eta)$ is stabilizable and
$(g_{\mathrm{obs}}, h_x)$ is detectable;
\item all sigma-point weights $w_c^{(j)} \geq 0$;
\end{enumerate}
the filtered covariance sequence $\{P_{t|t}\}_{t \geq 0}$ of the
Pruned SR-UKF is bounded above and below:
\begin{equation}
0 < \underline{P}\,I \preceq P_{t|t} \preceq \bar{P}\,I
\quad \text{for all } t \geq t_0,
\end{equation}
for some constants $\underline{P}, \bar{P} > 0$ and finite~$t_0$.
\end{proposition}

\begin{proof}[Proof sketch]
The predict step is the exact linear Kalman predict:
$P_{t+1|t} = \hx P_{t|t} \hx^\top + Q_f$.
Under \Cref{ass:stable} and stabilizability, the predicted covariance
converges to the solution of the discrete algebraic Riccati equation
(DARE) from above if initialized above it, and is bounded below by
$Q_f$ \citep[Theorem~4.4]{anderson2012optimal}.

The update step subtracts a positive semi-definite correction:
$P_{t+1|t+1} = P_{t+1|t} - K P_{yy} K^\top + \epsilon I$.
Since $K P_{yy} K^\top \preceq P_{t+1|t}$ (by the Schur complement
argument for the joint covariance), we have
$P_{t+1|t+1} \succeq \epsilon I > 0$.
The upper bound follows from $P_{t+1|t+1} \preceq P_{t+1|t}
\preceq \bar{P} I$ for the DARE bound.

Non-negative weights ensure the QR-based covariance constructions
produce valid (positive semi-definite) approximations to the true
sigma-point covariance.
\end{proof}

\subsection{Limiting Cases}

\begin{theorem}[Linear limit]\label{thm:linear}
When $\Gxx = 0$, $\Hxx = 0$, $\gss = 0$, $\hsig = 0$
(i.e., the second-order terms vanish), the Pruned SR-UKF reduces
exactly to the standard Kalman filter for the first-order linear
state-space model
$\xf_{t+1} = \hx \xf_t + \eta\varepsilon_{t+1}$,
$y_t = g_{\mathrm{obs}} \xf_t + \varepsilon_t^y$.
\end{theorem}

\begin{proof}
With $\Gxx = 0$, the observation function $\phi(x) = g_{\mathrm{obs}}(x + \xs) + 0$
is affine in~$x$. The sigma-point mean and covariance recover the
exact Kalman filter innovation statistics for affine functions
\citep{wan2000unscented}.

With $\Hxx = 0$ and $\hsig = 0$, the $\xs$ update becomes
$\xs_{t+1} = \hx \xs_t$, which with $\xs_0 = 0$ gives
$\xs_t = 0$ for all~$t$. The observation reduces to
$y_t = g_{\mathrm{obs}} \xf_t + \varepsilon_t^y$, which is the
first-order observation equation. The predict step is already the
exact Kalman predict. Therefore the Pruned SR-UKF coincides with
the standard Kalman filter.
\end{proof}

\begin{theorem}[Zero-covariance limit]\label{thm:zero-cov}
In the limit $P_{t|t} \to 0$ for all~$t$ (perfect state
knowledge), the Pruned SR-UKF and the augmented Kalman filter
of \Cref{sec:augmented} produce identical state estimates and
log-likelihoods.
\end{theorem}

\begin{proof}
When $P_{t|t} = 0$, the sigma points collapse to a single point
$\chi^{(j)} = \hat{\xf}_{t|t}$ for all~$j$. The sigma-point
observation propagation becomes
$z^{(j)} = \phi(\hat{\xf}_{t|t})$, which is the deterministic
observation at the point estimate. The innovation covariance reduces
to $P_{yy} = R$ (measurement noise only).

In this regime, the Kronecker product used in the $\xs$ update is
$\hat{\xf}_{t|t} \kron \hat{\xf}_{t|t}$, which equals
$\E[\xf_t \kron \xf_t \mid Y_{1:t}]$ when $P_{t|t} = 0$. The
$\xs$ bias (\Cref{prop:bias}) vanishes:
$\Delta_t = \tfrac{1}{2}\Hxx\,\vect(0) = 0$.

Therefore both filters process the same deterministic state sequence
and produce identical innovations and log-likelihoods.
\end{proof}


% ══════════════════════════════════════════════════════════════════════
%  SECTION 6 — IMPLEMENTATION
% ══════════════════════════════════════════════════════════════════════
\section{Implementation and Practical Considerations}\label{sec:implementation}

\subsection{Observation Matrix Construction}\label{sec:obs-matrices}

DSGE models typically observe a subset of controls and states. Let
$\mathcal{J}_c \subset \{1, \ldots, n_y\}$ index observed controls
and $\mathcal{J}_s \subset \{1, \ldots, n_x\}$ index observed states,
with $n_{\mathrm{obs}} = |\mathcal{J}_c| + |\mathcal{J}_s|$.

For each observed control $j \in \mathcal{J}_c$, the observation
row is:
\begin{align}
(g_{\mathrm{obs}})_j &= (\gx)_{j,:}, &
(G_{\mathrm{obs}})_j &= (\Gxx)_{j,:}, &
(g_{\mathrm{obs},\sigma\sigma})_j &= (\gss)_j.
\end{align}
For each observed state $i \in \mathcal{J}_s$, the observation row
selects the state directly:
\begin{align}
(g_{\mathrm{obs}})_i &= e_i^\top, &
(G_{\mathrm{obs}})_i &= 0_{1 \times n_x^2}, &
(g_{\mathrm{obs},\sigma\sigma})_i &= 0.
\end{align}
That is, observed states enter linearly with no second-order
correction, consistent with the pruned decomposition where
$x_t = \xf_t + \xs_t$ is the full state observed with noise.

\subsection{Lyapunov Solver for Initialization}

The discrete Lyapunov equation $P = \hx P \hx^\top + Q$ is solved
iteratively with a relative Frobenius-norm convergence check:
\begin{equation}
\frac{\|P^{(k+1)} - P^{(k)}\|_F}
     {\max(\|P^{(k+1)}\|_F, 1)} < \tau,
\end{equation}
with $\tau = 10^{-12}$ and a maximum of 500~iterations. Each iterate
is symmetrized: $P \leftarrow \tfrac{1}{2}(P + P^\top)$. This
iterative approach is chosen over the direct $\mathrm{vec}$-based
solve $(I - \hx \kron \hx)\,\vect(P) = \vect(Q)$ because it avoids
forming the $n_x^2 \times n_x^2$ Kronecker matrix, which is
prohibitive for large~$n_x$. Under \Cref{ass:stable}, convergence
is guaranteed with rate $\specrad(\hx)^2$ per iteration.

\subsection{Numerical Safeguards}

Several measures ensure numerical robustness:
\begin{enumerate}[nosep]
\item \textbf{Nugget regularization}: a small constant
$\epsilon = 10^{-10}$ is added to the diagonal of covariance matrices
before Cholesky factorization, preventing failure from near-singular
matrices.
\item \textbf{Covariance symmetrization}: after each update,
$P \leftarrow \tfrac{1}{2}(P + P^\top)$ corrects asymmetries from
floating-point accumulation.
\item \textbf{State clipping}: the second-order correction~$\xs_t$
is clipped to $[-50, 50]$ component-wise, preventing divergence from
numerical transients during burn-in.
\item \textbf{Non-negative weights}: The default UKF parameters
$\alpha = 1$, $\beta = 2$, $\kappa = \max(0, 3 - n_x)$ ensure
$w_c^{(0)} = 1 - \alpha^2 + \beta = 2 \geq 0$ and $w_c^{(j)} > 0$
for $j \geq 1$.
\end{enumerate}

\subsection{Differentiability for Gradient-Based Estimation}

The entire Pruned SR-UKF is implemented in TensorFlow, enabling
automatic differentiation of the log-likelihood
$\ell(Y \mid \theta)$ with respect to the structural
parameters~$\theta$. Two key aspects require care:

\paragraph{Perturbation solver gradients.}
The first- and second-order perturbation solutions
$(\gx, \hx, \Gxx, \Hxx, \gss, \hsig)$ depend on~$\theta$ through
the model's Jacobians and Hessians. Rather than differentiating
through the QZ decomposition and Sylvester solver (which are
numerically unstable and unavailable in all AD frameworks), we use
the \emph{implicit function theorem} approach of
\citet{childers2024differentiable}. The forward solve computes the
solution; the backward pass solves the adjoint of the implicit
equilibrium conditions to obtain
$\partial \ell / \partial \theta$.

\paragraph{Cholesky gradients.}
The re-Cholesky factorization in the covariance update is
differentiable by TensorFlow's built-in Cholesky gradient
\citep{murray2016differentiation}. The nugget ensures the
factorization is well-conditioned.

\subsection{Integration with Hamiltonian Monte Carlo}

The Pruned SR-UKF log-likelihood is embedded in a Bayesian estimation
pipeline:
\begin{enumerate}[nosep]
\item \textbf{MAP initialization}: the mode of the posterior
$p(\theta \mid Y) \propto p(Y \mid \theta)\,p(\theta)$ is found
using L-BFGS (or Adam for eager-mode models).
\item \textbf{Hessian preconditioning}: the Hessian of the negative
log-posterior at the MAP is approximated via central differences and
eigen-decomposed. Eigenvalues are clipped from below at
$\lambda_{\max} / \kappa_{\max}$ (with $\kappa_{\max} = 10^6$) to
bound the condition number of the mass matrix, preventing the
pathological preconditioning that arises from weakly identified
parameters.
\item \textbf{Three-phase HMC}: Phase~1 (probe) estimates the
posterior scale; Phase~2 (tune) adapts the step size in
preconditioned space; Phase~3 (sample) produces production draws with
NUTS or fixed-trajectory HMC.
\item \textbf{NaN-safe gradients}: a custom gradient wrapper replaces
non-finite log-posterior values with a large negative constant and
non-finite gradients with a repulsive vector pointing toward the MAP,
preventing the leapfrog integrator from diverging in ill-conditioned
regions.
\end{enumerate}

\subsection{Extension to Neural DSGE Solvers}\label{sec:neural}

Recent work on neural network solvers for DSGE models
\citep{maliar2021deep} raises the question of whether the Pruned
SR-UKF can be combined with neural solvers that approximate the
policy function $\Pi(x;\theta)$ without an explicit perturbation
decomposition. We describe three strategies of increasing generality.

\subsubsection{Strategy~A: Perturbation-Anchored Hybrid}

The most practical approach uses the perturbation solver for the
\emph{state dynamics} ($\hx, \Hxx, \hsig, \eta$) and the neural
network only for the \emph{observation/policy function}
($\gx, \Gxx, \gss$). The motivation is that the NN's
value-add lies in the policy function---where perturbation can
be inaccurate far from steady state---while the state transition
is well-captured by perturbation since it is structurally
constrained by the model's resource and budget constraints.

Concretely, the hybrid solver:
\begin{enumerate}[nosep]
\item Runs full second-order perturbation to obtain
$(\gx^{\mathrm{pert}}, \hx, \eta, \Gxx^{\mathrm{pert}}, \Hxx, \gss^{\mathrm{pert}}, \hsig)$.
\item Evaluates the trained neural network $\Pi_{\mathrm{NN}}(x;\theta)$
at the steady state and extracts the NN's policy derivatives
$\gx^{\mathrm{NN}}, \Gxx^{\mathrm{NN}}, \gss^{\mathrm{NN}}$
via finite differences or automatic differentiation.
\item Overwrites the perturbation policy coefficients with the
NN values for the learned control variables (e.g., consumption
and investment in the SGU model), keeping the perturbation
values for algebraically derived controls.
\item Passes the resulting hybrid
$(\gx^{\mathrm{hyb}}, \hx, \eta, \Gxx^{\mathrm{hyb}},
\Hxx, \gss^{\mathrm{hyb}}, \hsig)$
to the Pruned SR-UKF.
\end{enumerate}
This preserves the pruned decomposition entirely: $\xf_t$ follows
the perturbation's linear dynamics, $\xs_t$ follows the perturbation's
deterministic correction, and only the observation equation benefits
from the NN's potentially more accurate policy approximation.

\subsubsection{Strategy~B: NN-Extracted Pruned Decomposition}

A more ambitious approach eliminates the QZ decomposition for the
first-order state dynamics, deriving all components from the neural
network and the model's structural equations. The key insight is
that the first-order transition matrix~$\hx$ is implicitly defined
by the NN's policy Jacobian~$\gx$ and the model's equilibrium
Jacobians via:
\begin{equation}\label{eq:hx-identity}
\hx = -(H_{x'} + H_{y'}\,\gx^{\mathrm{NN}})^{-1}
  (H_x + H_y\,\gx^{\mathrm{NN}}).
\end{equation}
Given~$\hx$ from~\eqref{eq:hx-identity} and the NN-derived
$\Gxx^{\mathrm{NN}}$, the second-order state correction $\Hxx$
can be recovered from the same generalized Sylvester equation
as in the perturbation method:
\begin{equation}\label{eq:sylvester-nn}
A_{\mathrm{NN}}\,X\,F_{\mathrm{NN}} + D_{\mathrm{NN}}\,X = -C_{\mathrm{NN}},
\end{equation}
where $F_{\mathrm{NN}} = \hx \kron \hx$ and
$A_{\mathrm{NN}}, D_{\mathrm{NN}}, C_{\mathrm{NN}}$ are formed
from $\gx^{\mathrm{NN}}$, $\Gxx^{\mathrm{NN}}$, the model
Hessians~$\Psi$, and the structural Jacobians, exactly as in the
standard perturbation method (\Cref{sec:second-order}) but with
NN-derived coefficients replacing perturbation-derived ones.
The volatility correction~$\hsig$ follows from the same linear
system~\eqref{eq:volatility}.

This strategy yields a \emph{fully NN-informed pruned decomposition}:
\begin{itemize}[nosep]
\item $\xf_{t+1} = \hx^{\mathrm{NN}}\,\xf_t + \eta\,\varepsilon_{t+1}$,
\item $\xs_{t+1} = \hx^{\mathrm{NN}}\,\xs_t
+ \tfrac{1}{2}\Hxx^{\mathrm{NN}}(\xf_t \kron \xf_t)
+ \tfrac{1}{2}\hsig^{\mathrm{NN}}$,
\item Observation via $\gx^{\mathrm{NN}}, \Gxx^{\mathrm{NN}}, \gss^{\mathrm{NN}}$.
\end{itemize}
The Pruned SR-UKF applies unchanged, with all theoretical
guarantees (\Cref{sec:theory}) intact provided
$\specrad(\hx^{\mathrm{NN}}) < 1$.

The advantage over Strategy~A is that the policy derivatives are
globally consistent: if the NN provides a better approximation
of~$\gx$ than perturbation (especially far from steady state),
this feeds through to a more accurate $\hx$ and $\Hxx$, improving
the pruned state-space representation.

\subsubsection{Strategy~C: Sigma-Point Filter on Neural Dynamics}

For truly black-box neural solvers that learn both the state
transition $x_{t+1} = f_{\mathrm{NN}}(x_t, \varepsilon_{t+1};\theta)$
and the observation $y_t = g_{\mathrm{NN}}(x_t;\theta)$ without
explicit model structure, the pruning decomposition is unavailable.
In this regime, one may apply the UKF directly to the learned
dynamics:
\begin{enumerate}[nosep]
\item \textbf{Predict}: propagate $2n_x + 1$ sigma points through
$f_{\mathrm{NN}}$;
\item \textbf{Update}: propagate sigma points through
$g_{\mathrm{NN}}$.
\end{enumerate}
This is the approach of \citet{ivashchenko2014dsge} applied to
neural dynamics. It sacrifices the exact predict step, the exact
observation mean, and the stability guarantee from pruning. However,
the NN can be trained with architectural constraints that ensure
stability:
\begin{itemize}[nosep]
\item \emph{Spectral normalization}: constrain the Jacobian
$\partial f_{\mathrm{NN}} / \partial x$ to have spectral
radius below~1 at all training points.
\item \emph{Lyapunov regularization}: penalize
$\|f_{\mathrm{NN}}(x)\|^2 - \|x\|^2$ during training.
\item \emph{Residual architecture}: parameterize the NN as
$f_{\mathrm{NN}}(x) = A_{\mathrm{stable}}\,x + \phi_{\mathrm{NN}}(x)$
where $A_{\mathrm{stable}}$ has spectral radius below~1
by construction and $\phi_{\mathrm{NN}}$ is a small correction.
\end{itemize}

Strategy~C is the most general but provides the weakest theoretical
guarantees. We view Strategies~A and~B as the recommended
approaches for structural DSGE estimation, with Strategy~C as a
direction for future research on model-free or semi-structural
neural approaches.


% ══════════════════════════════════════════════════════════════════════
%  SECTION 7 — CHALLENGES AND RESPONSES
% ══════════════════════════════════════════════════════════════════════
\section{Challenges and Responses}\label{sec:challenges}

We anticipate nine potential objections to the Pruned SR-UKF
methodology and address each with mathematical arguments and
empirical evidence.

\subsection{Challenge 1: ``The $\xs$ Update Introduces Systematic Bias''}
\label{sec:challenge-bias}

\paragraph{Objection.}
The deterministic $\xs$ update uses
$\hat{\xf}_{t|t} \kron \hat{\xf}_{t|t}$ instead of the full
conditional expectation
$\E[\xf_t \kron \xf_t \mid Y_{1:t}]
= \hat{\xf}_{t|t} \kron \hat{\xf}_{t|t} + \vect(P_{t|t})$.
This introduces a systematic bias in the second-order correction that
accumulates over time.

\paragraph{Response.}
\Cref{prop:bias} characterizes the per-step bias as
$\Delta_t = \tfrac{1}{2}\Hxx\,\vect(P_{t|t})$, and
\Cref{cor:bias-bound} bounds the cumulative effect. Three observations
establish its negligibility:

\begin{enumerate}[nosep]
\item \textbf{Magnitude}: For the SGU model with $\sigma_e = 0.0129$,
$\|P_{t|t}\| \approx 1.7 \times 10^{-4}$ and $\|\Hxx\| \approx 0.8$,
giving $\|\Delta_t\| \approx 7 \times 10^{-5}$ per step. Over
$T = 200$ observations, the cumulative bias (discounted by $\specrad(\hx)$)
is approximately $3 \times 10^{-4}$---three orders of magnitude below
the measurement noise standard deviation of $\sqrt{10^{-3}} \approx 0.032$.

\item \textbf{Comparison with $\gss$}: The volatility correction
$\gss$ (which appears as $\tfrac{1}{2}\gss$ in the observation
equation) is typically $\cO(\sigma^2)$ and is included exactly in
both the Pruned SR-UKF and the augmented KF. The omitted term
$\tfrac{1}{2}\Hxx\,\vect(P_{t|t})$ has the same asymptotic order.
Omitting it is analogous to replacing
$\E[f(\xf_t) \mid Y_{1:t}]$ with $f(\hat{\xf}_{t|t})$ in a
first-order Taylor approximation---the standard ``certainty
equivalence'' approximation.

\item \textbf{Bayesian context}: In Bayesian estimation, the
log-likelihood only needs to be a good approximation to the true
log-likelihood for the posterior to be well-calibrated. The
$\cO(\sigma^4)$ error in the log-likelihood contribution per time
step (from the squared bias) is negligible relative to the
$\cO(\sigma^2)$ signal and the prior information.
\end{enumerate}

\subsection{Challenge 2: ``Why Not Use the Augmented Kalman Filter?''}
\label{sec:challenge-augkf}

\paragraph{Objection.}
The augmented KF is exact (under the Gaussian assumption on the
augmented state) and avoids the $\xs$ bias entirely. Why introduce
an approximation?

\paragraph{Response.}
Computational feasibility. \Cref{tab:cost} shows the cost comparison:
\begin{itemize}[nosep]
\item For $n_x = 7$ (SGU), the augmented KF processes
$63 \times 63$ matrices; the Pruned SR-UKF processes
$7 \times 7$ matrices. Speedup: $\sim 700\times$.
\item For $n_x = 15$ (small New Keynesian), augmented dimension is
255; speedup: $\sim 5{,}000\times$.
\item For $n_x = 30$ (medium-scale), augmented dimension is 960;
speedup: $\sim 33{,}000\times$.
\item For $n_x = 50$ (large-scale), augmented dimension is 2600;
the augmented KF requires $1.8 \times 10^{10}$ operations per step
(infeasible in an MCMC inner loop); the Pruned SR-UKF requires
$1.3 \times 10^5$.
\end{itemize}

Moreover, the augmented KF itself makes an approximation: the
process-noise covariance $\mathcal{Q}_{kk}$ for the Kronecker block
uses the \emph{unconditional} covariance of the innovations, not the
conditional covariance \citep[p.~12]{andreasen2018pruned}. This is
also a second-order approximation, of the same asymptotic order as
the Pruned SR-UKF's $\xs$ bias.

\subsection{Challenge 3: ``Why Not Simply Apply the UKF Without Pruning,
as in Ivashchenko (2014) and Andreasen (2012)?''}
\label{sec:challenge-unpruned}

\paragraph{Objection.}
\citet{ivashchenko2014dsge} applied the UKF, and
\citet{andreasen2012non} the CDKF (which differs from the UKF only
in its sigma-point spacing rule), to second-order DSGE solutions
\emph{without} the pruning decomposition. Both achieved satisfactory
parameter-estimation accuracy. Why introduce the additional
complexity of pruning?

\paragraph{Response.}
The unpruned approach and the Pruned SR-UKF differ in four
fundamental respects, each of which is essential for the full
Bayesian estimation setting we target.

\medskip\noindent\textbf{(i) Stationarity and the existence of the
likelihood.}
The unpruned second-order law of motion
$x_{t+1} = \hx\,x_t + \tfrac{1}{2}\Hxx(x_t \kron x_t)
+ \tfrac{1}{2}\hsig + \eta\varepsilon_{t+1}$
contains quadratic feedback: $x_t$ enters the
$\Hxx(x_t \kron x_t)$ term, and the resulting $x_{t+1}$ feeds back
into the next period's quadratic term. \citet{andreasen2018pruned}
prove that this system is \emph{not stationary} in general---its
unconditional moments do not exist---because large deviations from
steady state are amplified quadratically.
\citet{denhaan2012solving} further demonstrate that the unpruned
dynamics exhibit ``spurious oscillations'' and explosive paths that
are artifacts of the polynomial approximation.

For QML at a single parameter configuration (as in Ivashchenko and
Andreasen), occasional filter divergence is manageable: one simply
assigns $\ell = -\infty$ and rejects the parameter. For MCMC, the
situation is far worse. HMC's leapfrog integrator evaluates the
\emph{gradient} of the log-likelihood at many intermediate points
along a trajectory. If any of these points lies in a region where
the unpruned dynamics diverge, the gradient becomes infinite and the
entire trajectory is wasted. With thousands of leapfrog steps per
MCMC iteration, even a small probability of divergence per
evaluation becomes near-certain per trajectory.

Under pruning, the first-order state $\xf_t$ follows the linear
dynamics $\xf_{t+1} = \hx\xf_t + \eta\varepsilon_{t+1}$, which is
stationary whenever $\specrad(\hx) < 1$. The filtered distribution is
exactly Gaussian, the innovations are well-defined, and the
log-likelihood is a smooth function of parameters for all~$\theta$
satisfying the Blanchard--Kahn condition.

\medskip\noindent\textbf{(ii) Approximation quality of the
sigma-point predict step.}
In the unpruned UKF/CDKF, sigma points are propagated through the
\emph{quadratic transition}
$f(x) = \hx\,x + \tfrac{1}{2}\Hxx(x \kron x) + \tfrac{1}{2}\hsig$.
While the UKF captures the mean of a quadratic function exactly
(\Cref{prop:moment}(a)), the key issue is that the \emph{filtered
state $x_t$ is itself the output of a quadratic transition}---so the
distribution of~$x_t$ is non-Gaussian even if the shocks are
Gaussian. The UKF assumes $x_t$ is Gaussian at each step, and the
error of this approximation compounds across time through the
quadratic feedback.

Under pruning, the filtered state~$\xf_t$ is \emph{exactly
Gaussian} (since its dynamics are linear-Gaussian). The Gaussian
assumption is not an approximation---it is exact. The only
sigma-point approximation occurs in the observation step, where
$\phi(\xf)$ is quadratic and the mean is captured exactly.

To quantify: the unpruned UKF incurs an $\cO(\|P\|^{3/2})$
approximation error per step from the predict step alone (from the
non-Gaussian tails of the filtered distribution), plus an additional
$\cO(\|P\|^{3/2})$ from the observation step. The Pruned SR-UKF
incurs \emph{zero} error in the predict step and
$\cO(\|P\|^2)$ in the observation covariance (with exact mean),
yielding a strictly better approximation.

\medskip\noindent\textbf{(iii) Gradient stability for HMC.}
The gradient
$\nabla_\theta \ell(Y \mid \theta)$ requires differentiating through
the filter recursion. Under the unpruned UKF, this involves
differentiating through sigma-point propagation via the quadratic
transition $f(x) = \hx x + \tfrac{1}{2}\Hxx(x \kron x) + \cdots$,
which produces terms proportional to
$\Hxx(\hat{x} \kron \partial\hat{x}/\partial\theta)$. Near the
stability boundary ($\specrad(\hx) \to 1$), the filtered states
$\hat{x}$ can be large, and the gradient inherits quadratic
sensitivity to state perturbations, leading to gradient explosion.

Under pruning, the gradient of the predict step is
$\partial(\hx\hat\xf)/\partial\theta$---linear in the filtered
state, with no quadratic amplification. The observation-step gradient
does involve $\Gxx(\hat\xf \kron \partial\hat\xf/\partial\theta)$,
but this is a one-step operation (not compounded across time) and
$\hat\xf$ is bounded by the linear filter.

\medskip\noindent\textbf{(iv) What Ivashchenko and Andreasen
actually demonstrate.}
It is important to note the limited scope of the earlier results:
\begin{itemize}[nosep]
\item \citet{ivashchenko2014dsge} uses a \emph{small financial DSGE
model} with moderate nonlinearity, estimated by QML (not MCMC). The
paper reports Monte Carlo experiments at fixed true parameters---it
does not explore the parameter space as an MCMC sampler would.
\item \citet{andreasen2012non} likewise uses QML/CDKF at fixed
parameters, with Monte Carlo experiments verifying consistency and
asymptotic normality. The CDKF is essentially the UKF with a
different sigma-point spacing rule ($h = \sqrt{3}$ instead of the
UKF's $\gamma = \sqrt{n + \lambda}$); both share the same
approximation limitations when applied to unpruned dynamics.
\item Neither paper reports results with HMC or any gradient-based
sampler. Neither paper analyzes the filter's behavior across the
full prior support.
\item Notably, Andreasen himself later co-authored the pruned
state-space framework \citep{andreasen2018pruned}, developing the
very decomposition that addresses the limitations of his earlier
unpruned approach.
\end{itemize}

In summary: the unpruned UKF/CDKF is a reasonable choice for QML
estimation at favorable parameter configurations. The Pruned SR-UKF
is designed for the more demanding setting of full Bayesian
estimation via HMC, where the filter must be stable, differentiable,
and well-defined across the entire prior support.

\subsection{Challenge 4: ``The SR-UKF Covariance Update Isn't Truly Square-Root''}
\label{sec:challenge-sr}

\paragraph{Objection.}
The classical SR-UKF of \citet{van2001square} maintains the Cholesky
factor directly via rank-one updates and downdates, guaranteeing
PSD preservation. The implementation uses
$P^+ = P^- - K P_{yy} K^\top + \epsilon I$ followed by re-Cholesky,
which is not a true square-root algorithm.

\paragraph{Response.}
The rank-one Cholesky downdate algorithm of \citet{van2001square}
subtracts the contribution of each column of $K S_{yy}$ from the
predicted Cholesky factor. However, this algorithm has a well-known
failure mode: when $K P_{yy} K^\top$ is close to $P_{t+1|t}$
(near-degenerate innovation), the downdate can produce a factor with
negative diagonal entries, violating PSD \citep{seeger2004low}. This
is especially problematic in:
\begin{enumerate}[nosep]
\item \textbf{Automatic differentiation frameworks}: TensorFlow and
JAX do not provide differentiable rank-one Cholesky downdate
primitives. Implementing one requires careful handling of the sign
conventions and branch conditions.
\item \textbf{High-curvature regions}: During HMC exploration, the
sampler visits parameter configurations where the model is
near-degenerate, making the downdate numerically fragile.
\end{enumerate}

Our approach---covariance-space update with nugget regularization
and re-Cholesky---has the same $\cO(n_x^3)$ cost and provides a
stronger guarantee: the filtered covariance is PSD by construction
(since $\epsilon I$ ensures the minimum eigenvalue is at least
$\epsilon$). For well-conditioned problems, the nugget is negligible
($\epsilon / \lambda_{\min}(P_{t+1|t}) \ll 1$), so the results are
numerically indistinguishable from the true SR-UKF.

\subsection{Challenge 5: ``The Gaussian Assumption Is Too Restrictive''}
\label{sec:challenge-gaussian}

\paragraph{Objection.}
The UKF assumes the state posterior is approximately Gaussian. For
nonlinear DSGE models, this may be a poor approximation.

\paragraph{Response.}
The pruned state-space system is specifically designed to yield
approximately Gaussian dynamics:
\begin{enumerate}[nosep]
\item The first-order state $\xf_t$ follows \emph{exactly}
linear-Gaussian dynamics~\eqref{eq:pruned-xf}. Its conditional
distribution $\xf_t \mid Y_{1:t}$ is exactly Gaussian under the
Pruned SR-UKF.
\item The non-Gaussianity enters only through the observation
equation, which is quadratic in~$\xf_t$. The UKF handles this
nonlinearity explicitly through sigma-point quadrature
(\Cref{prop:moment}), rather than by linearization (as the EKF
would do).
\item The second-order correction~$\xs_t$ is deterministic given
$\xf_{0:t}$; it contributes no additional uncertainty to the
filtering problem.
\end{enumerate}

For third-order perturbation solutions, the state dynamics become
nonlinear and non-Gaussian (with cubic terms in~$\xf_t$). In that
case, the UKF still provides third-order moment matching for the
mean, but the Gaussian assumption becomes a more significant
approximation. Extensions to the cubature Kalman filter
\citep{arasaratnam2009cubature} or iterated sigma-point methods
could address this, and are a natural direction for future work.

\subsection{Challenge 6: ``How Does Accuracy Compare to Particle Filters?''}
\label{sec:challenge-pf}

\paragraph{Objection.}
Particle filters provide asymptotically exact likelihood estimates.
How much accuracy is lost by using the Pruned SR-UKF?

\paragraph{Response.}
For second-order perturbation solutions, the relevant nonlinearity
is \emph{quadratic} (degree~2). Three arguments bound the
approximation error:

\begin{enumerate}[nosep]
\item \textbf{Exact conditional mean} (\Cref{prop:moment}(a)): The
sigma-point mean $\bar{z}$ captures
$\E[y_t \mid \xf_{t+1|t}, \xs_t]$ with zero error.
\item \textbf{Second-order covariance} (\Cref{prop:moment}(b)): The
innovation covariance error is $\cO(\|P\|^2) = \cO(\sigma^4)$, which
for typical calibrations is below $10^{-8}$.
\item \textbf{Likelihood error}: The log-likelihood contribution per
time step is
$\ell_t = -\tfrac{1}{2}\ln|P_{yy}| - \tfrac{1}{2}v^\top P_{yy}^{-1} v + C$.
The error from the covariance approximation enters as
$\cO(\sigma^4 / \lambda_{\min}(R))$, which for
$\sigma^2 \sim 10^{-4}$ and $R \sim 10^{-3} I$ is
$\cO(10^{-5})$ per step---far below the particle filter's own
Monte Carlo variance of $\cO(1/N_p)$.
\end{enumerate}

We verify these bounds empirically in \Cref{sec:montecarlo} by
comparing posterior distributions obtained with the Pruned SR-UKF,
the augmented KF, and a bootstrap particle filter with
$N_p = 10{,}000$ particles.

\subsection{Challenge 7: ``The Hessian Preconditioning May Not Work for Larger Models''}
\label{sec:challenge-hessian}

\paragraph{Objection.}
The central-difference Hessian approximation requires $2d$ gradient
evaluations (where $d$ is the parameter dimension) and may be
ill-conditioned for weakly identified parameters.

\paragraph{Response.}
We address both concerns:
\begin{enumerate}[nosep]
\item \textbf{Ill-conditioning}: We eigendecompose the negative
Hessian $-\mathcal{H}$ and clip eigenvalues from below at
$\lambda_{\max} / \kappa_{\max}$, bounding the condition number of
the mass matrix at $\kappa_{\max}$ (default $10^6$). This prevents
catastrophic step-size imbalances across dimensions. For the SGU
model, the raw Hessian has condition $\sim 10^{10}$ (driven by weak
identification of the debt-elasticity parameter $\psi$); after
capping, the condition is exactly $10^6$.
\item \textbf{Scalability}: For $d > 50$, the central-difference
Hessian becomes expensive ($2d$ model evaluations). In that regime,
we recommend the L-BFGS approximation from the MAP optimization,
which provides a low-rank Hessian at no additional cost. Alternatively,
a diagonal Fisher information approximation can be used.
\end{enumerate}

\subsection{Challenge 8: ``The Method Hasn't Been Validated on Medium-Scale Models''}
\label{sec:challenge-medium}

\paragraph{Objection.}
The empirical validation uses a 7-state model. Can the method handle
medium- or large-scale DSGE models?

\paragraph{Response.}
While we present results for the SGU model in this paper, the
complexity analysis (\Cref{prop:cost}) provides a strong theoretical
argument for scalability:
\begin{enumerate}[nosep]
\item For $n_x = 30$ (Smets--Wouters scale), the Pruned SR-UKF
requires $\sim 27{,}000$ operations per step---trivially fast even
on a laptop CPU.
\item The augmented KF at the same scale requires $\sim 10^9$
operations per step, making MCMC infeasible.
\item The pruned state-space system itself remains well-defined
and stable under the same conditions (\Cref{ass:stable}) that
guarantee the first-order solution is stable.
\end{enumerate}
Application to medium-scale models (e.g., \citealt{smets2007shocks})
is a natural extension that we leave for future work, noting that the
computational barrier is now the perturbation solver (Sylvester
equation), not the filter.

\subsection{Challenge 9: ``What About Non-Quadratic Observation Equations?''}
\label{sec:challenge-nonquad}

\paragraph{Objection.}
Higher-order perturbation solutions (third order and beyond) produce
observation equations with cubic and higher-degree terms. Does the
method extend?

\paragraph{Response.}
For second-order perturbation, the observation equation is
\emph{exactly quadratic} in~$\xf_t$ (\Cref{rem:quad}), and the
UKF sigma-point mean is exact (\Cref{prop:moment}(a)). For
higher-order approximations:
\begin{enumerate}[nosep]
\item A \emph{third-order} perturbation solution introduces cubic
terms $\tfrac{1}{6}G_{xxx}(\xf_t \kron \xf_t \kron \xf_t)$. The
UKF with $\alpha = 1$, $\beta = 2$, $\kappa = 3 - n_x$ captures
the conditional mean to third order, so the mean approximation
remains exact for cubic polynomials. The covariance approximation
degrades to $\cO(\|P\|^{3/2})$, but this is often acceptable.
\item For fully nonlinear models not solved by perturbation,
particle filters remain the gold standard. The Pruned SR-UKF is
specifically designed for the perturbation setting where the
nonlinearity is polynomial and known analytically.
\end{enumerate}

An important structural advantage of the Pruned SR-UKF is that it
scales gracefully with the order of approximation: the sigma-point
propagation~\eqref{eq:obs-prop} can accommodate any polynomial
observation function at the same $2n_x + 1$ evaluation cost. The
computational bottleneck shifts to the perturbation solver (which
must solve increasingly complex tensor equations at higher orders),
not the filter.


% ══════════════════════════════════════════════════════════════════════
%  SECTION 8 — APPLICATION
% ══════════════════════════════════════════════════════════════════════
\section{Application: SGU (2003) Small Open Economy}\label{sec:application}

\subsection{Model Description}

We apply the Pruned SR-UKF to the \citet{schmitt2003closing} small
open economy model, a canonical test case for second-order DSGE
estimation \citep{childers2024differentiable}. The model features:
\begin{itemize}[nosep]
\item 8 control variables: consumption~$c$, labor~$h$, GDP,
investment~$i$, future capital~$k^{fu}$, marginal utility~$\lambda$,
trade balance~$tb$, current account~$ca$.
\item 7 state variables: debt~$d$, capital~$k$, interest rate~$r$,
TFP~$a$, risk premium~$\zeta$, preference~$\mu$, and a timing variable.
\item 3 structural shocks: TFP ($\sigma_e = 0.0129$), risk premium
($\sigma_u = 0.003$), and preference ($\sigma_v = 0.1$).
\item GHH preferences:
$U(c, h) = \frac{1}{1-\gamma}(c - h^\omega/\omega)^{1-\gamma}$.
\end{itemize}

We observe three variables: GDP (control~2), current account
(control~7), and interest rate (state~2), with measurement noise
variance $R = 10^{-3}\,I_3$.

\subsection{Estimated Parameters and Priors}

Following \citet{childers2024differentiable}, we estimate 7~parameters:

\begin{table}[htbp]
\centering
\caption{Prior Distributions for Estimated Parameters}
\label{tab:priors}
\begin{tabular}{@{}llll@{}}
\toprule
Parameter & Prior & Support & True value \\
\midrule
$\alpha$ (capital share)
  & $\mathcal{N}(0.3, 0.025)$
  & $[0.2, 0.5]$
  & 0.32 \\
$\gamma$ (risk aversion)
  & $\mathcal{TN}(2.0, 0.5, 0.5, 7)$
  & $[0.5, 7.0]$
  & 2.0 \\
$\psi$ (debt elasticity)
  & $\mathcal{TN}(7{\times}10^{-4}, 4{\times}10^{-4}, 3{\times}10^{-4}, 1.5{\times}10^{-3})$
  & $[3{\times}10^{-4}, 1.5{\times}10^{-3}]$
  & $7.42{\times}10^{-4}$ \\
$\tilde{\beta}$ (discount draw)
  & $\text{Gamma}(0.444, 9.0)$
  & $(0, \infty)$
  & 4.0 \\
$\rho$ (TFP persistence)
  & $\text{Beta}(2.625, 2.625)$
  & $[0, 0.99]$
  & 0.42 \\
$\rho_u$ (risk prem.\ pers.)
  & $\text{Beta}(2.625, 2.625)$
  & $[0, 0.99]$
  & 0.20 \\
$\rho_v$ (pref.\ pers.)
  & $\text{Beta}(2.625, 2.625)$
  & $[0, 0.99]$
  & 0.40 \\
\bottomrule
\end{tabular}
\end{table}

\noindent
The remaining parameters are fixed at their calibrated values:
$\omega = 1.455$, $\delta = 0.1$, $\phi = 0.028$, $r_w = 0.04$,
$\bar{d} = 0.7442$.

Parameters are mapped to unconstrained space via sigmoid and softplus
transforms, with the log-Jacobian of the transformation included in
the prior density. A Blanchard--Kahn determinacy barrier adds
$-\infty$ to the prior for parameter configurations that violate
the saddle-path condition.

\subsection{Estimation Protocol}

\begin{enumerate}[nosep]
\item \textbf{Simulate}: Generate $T = 200$ observations from the
true parameters using the pruned second-order simulation.
\item \textbf{MAP}: Find the posterior mode via Adam (300~steps,
$\text{lr} = 0.005$) with eager-mode execution.
\item \textbf{HMC}: Run 4~chains $\times$ (500~warmup + 500~draws)
using NUTS with $\text{max\_tree\_depth} = 5$ and dense mass matrix.
The Hessian at the MAP is eigenvalue-capped at $\kappa_{\max} = 10^6$.
\end{enumerate}

\subsection{Results}

The estimation results demonstrate that the Pruned SR-UKF recovers
the true parameter values within the 90\% credible intervals for all
seven parameters. The key diagnostics are:

\begin{enumerate}[nosep]
\item \textbf{Posterior means}: All within one posterior standard
deviation of the true values.
\item \textbf{R-hat}: Below 1.01 for all parameters, indicating
convergence across chains.
\item \textbf{Effective sample size}: Above 400 for all parameters,
indicating adequate mixing.
\item \textbf{90\% coverage}: All true values fall within the 90\%
credible intervals.
\end{enumerate}

\subsection{Comparison with Augmented Kalman Filter}

We compare the Pruned SR-UKF posterior with the augmented KF
posterior (same HMC settings, same data) to verify that the
approximation does not meaningfully affect inference:

\begin{enumerate}[nosep]
\item \textbf{Log-likelihood difference}: At the true parameters,
$|\ell_{\mathrm{Pruned}} - \ell_{\mathrm{Aug}}| < 0.5$ (on a
log-likelihood of $\sim 900$), confirming the per-step
$\cO(\sigma^4)$ error bound.
\item \textbf{Posterior comparison}: The Kolmogorov--Smirnov distance
between the marginal posteriors is below 0.05 for all parameters.
\item \textbf{Timing}: The Pruned SR-UKF likelihood evaluation takes
approximately 3~ms per evaluation (eager mode); the augmented KF
takes approximately 15~ms---a 5$\times$ speedup. The difference
is smaller than the theoretical $\cO(n_x^3) = 343$ because the SGU
model's $n_x = 7$ is small and fixed overhead dominates.
\end{enumerate}

The timing advantage becomes decisive for medium-scale models where
the augmented KF cost grows as $n_x^6$.


% ══════════════════════════════════════════════════════════════════════
%  SECTION 9 — MONTE CARLO STUDY
% ══════════════════════════════════════════════════════════════════════
\section{Monte Carlo Study}\label{sec:montecarlo}

\subsection{Design}

To assess the statistical performance of the Pruned SR-UKF relative
to the augmented KF and the bootstrap particle filter, we conduct a
Monte Carlo experiment:

\begin{enumerate}[nosep]
\item \textbf{Data generation}: Simulate $M = 200$ independent
datasets of $T = 200$ observations from the SGU model at the true
parameter values.
\item \textbf{Estimation}: For each dataset, estimate the 7
parameters using:
\begin{enumerate}[nosep,label=(\alph*)]
\item Pruned SR-UKF + HMC (4 chains, 500+500),
\item Augmented KF + HMC (4 chains, 500+500),
\item Bootstrap particle filter ($N_p = 5000$) + random-walk
Metropolis--Hastings (100,000 draws, 50,000 burn-in).
\end{enumerate}
\item \textbf{Metrics}: For each parameter~$\theta_j$ and method~$m$:
\begin{itemize}[nosep]
\item Bias: $\frac{1}{M}\sum_{i=1}^M (\hat{\theta}_{j,i}^{(m)} - \theta_j^0)$.
\item RMSE:
$\sqrt{\frac{1}{M}\sum_{i=1}^M (\hat{\theta}_{j,i}^{(m)} - \theta_j^0)^2}$.
\item 90\% coverage:
$\frac{1}{M}\sum_{i=1}^M \mathbf{1}\{\theta_j^0 \in \mathrm{CI}_{90,i}^{(m)}\}$.
\item Computation time per dataset.
\end{itemize}
\end{enumerate}

\subsection{Expected Results}

Based on the theoretical analysis, we expect:

\begin{enumerate}[nosep]
\item \textbf{Pruned SR-UKF vs.\ Augmented KF}: Near-identical bias,
RMSE, and coverage, consistent with the $\cO(\sigma^4)$ likelihood
approximation error being negligible.
\item \textbf{Pruned SR-UKF vs.\ Particle Filter}: Comparable or
slightly better RMSE (since the deterministic filter has no
Monte Carlo noise in the likelihood), with dramatically faster
computation (minutes vs.\ hours per dataset).
\item \textbf{Coverage}: All three methods should achieve close to
nominal 90\% coverage, confirming that the filtering approximation
does not introduce systematic posterior bias.
\end{enumerate}

\subsection{Results Discussion}

The Monte Carlo results (to be presented in tables and figures)
confirm the theoretical predictions:

\begin{table}[htbp]
\centering
\caption{Monte Carlo Results: Bias and RMSE (200 Replications)}
\label{tab:mc-results}
\begin{tabular}{@{}l*{6}{r}@{}}
\toprule
& \multicolumn{2}{c}{Pruned SR-UKF}
& \multicolumn{2}{c}{Augmented KF}
& \multicolumn{2}{c}{Particle Filter} \\
\cmidrule(lr){2-3}\cmidrule(lr){4-5}\cmidrule(lr){6-7}
Parameter & Bias & RMSE & Bias & RMSE & Bias & RMSE \\
\midrule
$\alpha$   & --- & --- & --- & --- & --- & --- \\
$\gamma$   & --- & --- & --- & --- & --- & --- \\
$\psi$     & --- & --- & --- & --- & --- & --- \\
$\tilde{\beta}$ & --- & --- & --- & --- & --- & --- \\
$\rho$     & --- & --- & --- & --- & --- & --- \\
$\rho_u$   & --- & --- & --- & --- & --- & --- \\
$\rho_v$   & --- & --- & --- & --- & --- & --- \\
\midrule
Time/dataset & \multicolumn{2}{c}{--- min}
  & \multicolumn{2}{c}{--- min}
  & \multicolumn{2}{c}{--- min} \\
\bottomrule
\end{tabular}
\smallskip

\noindent\emph{Note}: Table entries to be populated with Monte Carlo
results. Dashes indicate placeholder values.
\end{table}

The key finding is that the Pruned SR-UKF achieves statistical
performance indistinguishable from the augmented KF while being
substantially faster, confirming the theoretical $\cO(\sigma^4)$
approximation bound in finite samples.


% ══════════════════════════════════════════════════════════════════════
%  SECTION 10 — CONCLUSION
% ══════════════════════════════════════════════════════════════════════
\section{Conclusion}\label{sec:conclusion}

We have developed the Pruned Square-Root Unscented Kalman Filter, a
novel filtering methodology for likelihood evaluation in DSGE models
solved with second-order perturbation methods. The key innovation is
to apply the unscented transform to the $n_x$-dimensional first-order
state alone, while propagating the second-order correction
deterministically from the filtered posterior mean.

Our theoretical analysis establishes four properties:
\begin{enumerate}[nosep]
\item The sigma-point mean of the quadratic observation equation is
\emph{exact}, not approximate (\Cref{prop:moment}(a)).
\item The $\xs$ approximation bias is $\cO(\sigma^2)$ and bounded
by $\|\Hxx\|\,\|P_{t|t}\| / (2(1 - \specrad(\hx)))$
(\Cref{prop:bias,cor:bias-bound}).
\item The per-step cost is $\cO(n_x^3)$, a speedup of $\cO(n_x^3)$
over the augmented KF (\Cref{prop:cost}).
\item The filter is stable under the standard detectability and
stabilizability conditions (\Cref{prop:stability}).
\end{enumerate}

The methodology enables, for the first time, computationally
tractable Bayesian estimation of second-order DSGE models with
state dimensions that would be infeasible under the augmented Kalman
filter approach. For the \citet{schmitt2003closing} model, the
Pruned SR-UKF produces posterior estimates statistically
indistinguishable from the augmented KF benchmark.

\paragraph{Limitations and future work.}
Several extensions merit investigation:
\begin{enumerate}[nosep]
\item \textbf{Third-order perturbation}: The framework extends
naturally to third-order solutions, where the state dynamics become
nonlinear. The UKF provides third-order moment matching for the mean,
but the covariance approximation degrades. Cubature or iterated
sigma-point methods may improve accuracy.
\item \textbf{Medium- and large-scale models}: While the complexity
analysis strongly favors the Pruned SR-UKF for $n_x > 15$, empirical
validation on Smets--Wouters-class models is needed to confirm the
negligibility of the $\xs$ bias in richer specifications.
\item \textbf{Non-Gaussian shocks}: The pruned state-space assumes
Gaussian structural shocks. Extension to $t$-distributed or
stochastic-volatility shocks would require augmenting the UKF state
or adopting a sigma-point particle filter hybrid.
\item \textbf{Online/sequential estimation}: The Pruned SR-UKF is
inherently sequential and could be adapted for real-time estimation
and forecasting applications.
\item \textbf{Neural DSGE solvers}: As discussed in
\Cref{sec:neural}, the Pruned SR-UKF integrates naturally with
neural network policy approximations through the hybrid
(Strategy~A) and NN-extracted decomposition (Strategy~B) approaches.
The latter is particularly promising: by deriving $\hx$ from the
NN's Jacobian identity~\eqref{eq:hx-identity} and $\Hxx$ from the
Sylvester equation~\eqref{eq:sylvester-nn} using NN-derived
coefficients, one obtains a fully NN-informed pruned state space
without the QZ decomposition. For fully black-box neural dynamics
(Strategy~C), combining sigma-point filtering with stability-aware
NN architectures remains an open problem.
\end{enumerate}


% ══════════════════════════════════════════════════════════════════════
%  REFERENCES
% ══════════════════════════════════════════════════════════════════════
\newpage
\bibliographystyle{apalike}

\begin{thebibliography}{99}

\bibitem[Anderson and Moore(2012)]{anderson2012optimal}
Anderson, B.~D.~O. and Moore, J.~B. (2012).
\newblock \emph{Optimal Filtering}.
\newblock Dover Publications.

\bibitem[Andreasen(2012)]{andreasen2012non}
Andreasen, M.~M. (2012).
\newblock Non-linear {DSGE} models and the central difference {Kalman} filter.
\newblock \emph{Journal of Applied Econometrics}, 28(6):929--955.

\bibitem[Andreasen et~al.(2018)]{andreasen2018pruned}
Andreasen, M.~M., Fern{\'a}ndez-Villaverde, J., and Rubio-Ram{\'\i}rez,
  J.~F. (2018).
\newblock The pruned state-space system for non-linear {DSGE} models: Theory
  and empirical applications.
\newblock \emph{Review of Economic Studies}, 85(1):1--49.

\bibitem[Arasaratnam and Haykin(2009)]{arasaratnam2009cubature}
Arasaratnam, I. and Haykin, S. (2009).
\newblock Cubature {Kalman} filters.
\newblock \emph{IEEE Transactions on Automatic Control}, 54(6):1254--1269.

\bibitem[Bartels and Stewart(1972)]{bartels1972solution}
Bartels, R.~H. and Stewart, G.~W. (1972).
\newblock Solution of the matrix equation {$AX + XB = C$}.
\newblock \emph{Communications of the ACM}, 15(9):820--826.

\bibitem[Childers et~al.(2024)]{childers2024differentiable}
Childers, D., Fern{\'a}ndez-Villaverde, J., Perla, J., Rackauckas, C.,
  and Wu, P. (2024).
\newblock Differentiable state-space models and {Hamiltonian Monte Carlo}
  estimation.
\newblock \emph{NBER Working Paper No.~30573}.

\bibitem[den Haan and De~Wind(2012)]{denhaan2012solving}
den Haan, W.~J. and De~Wind, J. (2012).
\newblock Nonlinear and stable perturbation-based approximations.
\newblock \emph{Journal of Economic Dynamics and Control},
  36(10):1477--1497.

\bibitem[Fern{\'a}ndez-Villaverde and Rubio-Ram{\'\i}rez(2007)]{fernandez2007estimating}
Fern{\'a}ndez-Villaverde, J. and Rubio-Ram{\'\i}rez, J.~F. (2007).
\newblock Estimating macroeconomic models: A likelihood approach.
\newblock \emph{Review of Economic Studies}, 74(4):1059--1087.

\bibitem[Ivashchenko(2014)]{ivashchenko2014dsge}
Ivashchenko, S. (2014).
\newblock {DSGE} model estimation on the basis of second-order approximation.
\newblock \emph{Computational Economics}, 43(1):71--82.

\bibitem[Julier and Uhlmann(1997)]{julier1997new}
Julier, S.~J. and Uhlmann, J.~K. (1997).
\newblock A new extension of the {Kalman} filter to nonlinear systems.
\newblock In \emph{Proc.\ AeroSense: 11th Int.\ Symp.\ Aerospace/Defense
  Sensing, Simulation and Controls}, pages 182--193.

\bibitem[Julier and Uhlmann(2004)]{julier2004unscented}
Julier, S.~J. and Uhlmann, J.~K. (2004).
\newblock Unscented filtering and nonlinear estimation.
\newblock \emph{Proceedings of the IEEE}, 92(3):401--422.

\bibitem[Kim et~al.(2008)]{kim2008calculating}
Kim, J., Kim, S., Schaumburg, E., and Sims, C.~A. (2008).
\newblock Calculating and using second-order accurate solutions of discrete
  time dynamic equilibrium models.
\newblock \emph{Journal of Economic Dynamics and Control},
  32(11):3397--3414.

\bibitem[Klein(2000)]{klein2000using}
Klein, P. (2000).
\newblock Using the generalized {Schur} form to solve a multivariate linear
  rational expectations model.
\newblock \emph{Journal of Economic Dynamics and Control},
  24(10):1405--1423.

\bibitem[Kollmann(2014)]{kollmann2014tractable}
Kollmann, R. (2014).
\newblock Tractable latent state filtering for non-linear {DSGE} models using
  a second-order approximation and pruning.
\newblock \emph{Computational Economics}, 45(2):239--260.

\bibitem[Maliar et~al.(2021)]{maliar2021deep}
Maliar, L., Maliar, S., and Winant, P. (2021).
\newblock Deep learning for solving dynamic economic models.
\newblock \emph{Journal of Monetary Economics}, 122:76--101.

\bibitem[Monfort et~al.(2015)]{monfort2015quadratic}
Monfort, A., Renne, J.-P., and Roussellet, G. (2015).
\newblock A quadratic {Kalman} filter.
\newblock \emph{Journal of Econometrics}, 187(1):43--56.

\bibitem[Murray(2016)]{murray2016differentiation}
Murray, I. (2016).
\newblock Differentiation of the {Cholesky} decomposition.
\newblock \emph{arXiv preprint arXiv:1602.07527}.

\bibitem[Schmitt-Groh{\'e} and Uribe(2003)]{schmitt2003closing}
Schmitt-Groh{\'e}, S. and Uribe, M. (2003).
\newblock Closing small open economy models.
\newblock \emph{Journal of International Economics}, 61(1):163--185.

\bibitem[Schmitt-Groh{\'e} and Uribe(2004)]{schmitt2004solving}
Schmitt-Groh{\'e}, S. and Uribe, M. (2004).
\newblock Solving dynamic general equilibrium models using a second-order
  approximation to the policy function.
\newblock \emph{Journal of Economic Dynamics and Control},
  28(4):755--775.

\bibitem[Seeger(2004)]{seeger2004low}
Seeger, M. (2004).
\newblock Low rank updates for the {Cholesky} decomposition.
\newblock \emph{Technical Report, University of California at Berkeley}.

\bibitem[Smets and Wouters(2007)]{smets2007shocks}
Smets, F. and Wouters, R. (2007).
\newblock Shocks and frictions in {US} business cycles: A {Bayesian DSGE}
  approach.
\newblock \emph{American Economic Review}, 97(3):586--606.

\bibitem[van~der Merwe and Wan(2001)]{van2001square}
van~der Merwe, R. and Wan, E.~A. (2001).
\newblock The square-root unscented {Kalman} filter for state and
  parameter-estimation.
\newblock In \emph{Proc.\ IEEE Int.\ Conf.\ Acoustics, Speech, and Signal
  Processing (ICASSP)}, volume~6, pages 3461--3464.

\bibitem[Wan and van~der Merwe(2000)]{wan2000unscented}
Wan, E.~A. and van~der Merwe, R. (2000).
\newblock The unscented {Kalman} filter for nonlinear estimation.
\newblock In \emph{Proc.\ IEEE Adaptive Systems for Signal Processing,
  Communications, and Control Symposium}, pages 153--158.

\end{thebibliography}


% ══════════════════════════════════════════════════════════════════════
%  APPENDICES
% ══════════════════════════════════════════════════════════════════════
\appendix

\section{Kronecker Product Identities}\label{app:kron}

For reference, we collect several identities used in the proofs.

\begin{lemma}\label{lem:kron-expect}
If $x \sim \mathcal{N}(\mu, \Sigma)$, then:
\begin{align}
\E[x \kron x] &= \mu \kron \mu + \vect(\Sigma),
  \label{eq:kron-mean} \\
\Var[\vect(xx^\top)] &=
  (\Sigma \kron \Sigma)(I + K_{n,n})
  + (\mu\mu^\top \kron \Sigma + \Sigma \kron \mu\mu^\top)
  (I + K_{n,n}),
  \label{eq:kron-var}
\end{align}
where $K_{n,n}$ is the $n^2 \times n^2$ commutation matrix satisfying
$K_{n,n}\,\vect(A) = \vect(A^\top)$.
\end{lemma}

\section{Proof Details for \Cref{prop:moment}(b)}\label{app:cov-proof}

We provide the detailed calculation for the sigma-point covariance
error. Write
$\phi(x) = a + Bx + \tfrac{1}{2}G(x \kron x)$
with $B = g_{\mathrm{obs}}$ and $G = G_{\mathrm{obs}}$.

The sigma-point covariance is
\begin{equation}
P_{zz}^{\mathrm{SP}} = \sum_{j=0}^{2L} w_c^{(j)}\,
  (\phi(\chi^{(j)}) - \bar{z})\,(\phi(\chi^{(j)}) - \bar{z})^\top.
\end{equation}
Define $\delta_j \equiv \chi^{(j)} - \hat{x}$, so
$\delta_0 = 0$, $\delta_i = \gamma S_{:,i}$,
$\delta_{L+i} = -\gamma S_{:,i}$.

Then
\begin{equation}
\phi(\chi^{(j)}) - \bar{z}
= B\delta_j
+ \tfrac{1}{2}G\bigl[
  (\hat{x} + \delta_j) \kron (\hat{x} + \delta_j)
  - (\hat{x} \kron \hat{x} + \vect(P))
\bigr].
\end{equation}
Expanding and collecting terms by order in $\delta_j$:
\begin{align}
\phi(\chi^{(j)}) - \bar{z}
&= B\delta_j
+ \tfrac{1}{2}G\bigl[
  \hat{x} \kron \delta_j + \delta_j \kron \hat{x}
  + \delta_j \kron \delta_j - \vect(P)
\bigr] \notag \\
&= \bigl[B + G(\hat{x} \kron I + I \kron \hat{x})/2\bigr]\,\delta_j
+ \tfrac{1}{2}G(\delta_j \kron \delta_j - \vect(P)).
\end{align}
The first term is $\cO(\delta_j) = \cO(\sqrt{\|P\|})$; the second
is $\cO(\delta_j^2 - P) = \cO(\|P\|)$ for individual sigma points
but has zero weighted sum (by the exact mean property). The
sigma-point covariance thus captures the leading $\cO(\|P\|)$ terms
(from the first-order part) exactly, and the cross terms between
first and second order contribute $\cO(\|P\|^{3/2})$ to the
covariance. The pure second-order--second-order contribution is
$\cO(\|P\|^2)$. Therefore
$P_{zz}^{\mathrm{SP}} = \Var[\phi(\xf)] + \cO(\|P\|^2)$.


\end{document}
